% !TEX program = xelatex
\documentclass[12pt,a4paper]{article}

% ── Packages ─────────────────────────────────────────────────────
\usepackage[utf8]{inputenc}
\usepackage[T1]{fontenc}
\usepackage[numbers]{natbib}
\usepackage[english]{babel}
\usepackage{amsmath, amssymb, amsthm}
\usepackage{mathtools}
\usepackage{geometry}
\usepackage{booktabs}
\usepackage{array}
\usepackage{longtable}
\usepackage{graphicx}
\usepackage{xcolor}
\usepackage{hyperref}
\usepackage{setspace}
\usepackage{titlesec}
\usepackage{abstract}
\usepackage{enumitem}
\usepackage{fancyhdr}
\usepackage{float}
\usepackage{caption}
\usepackage{subcaption}
\usepackage{tcolorbox}
\usepackage{mdframed}
\usepackage{times}
\usepackage{iftex}
\usepackage{fontspec}
\usepackage{tikz}
\usetikzlibrary{calc}

\newtcolorbox{keyeq}{
	enhanced, breakable,
	colback=white, colframe=fxgray,
	boxrule=0.6pt, arc=2pt,
	left=6pt, right=6pt, top=4pt, bottom=4pt
}

% ── Path to figures ───────────────────────────────────────────
% Extract Graphs.rar into a subfolder "graphs/" placed next to
% this .tex file before compiling.
\graphicspath{{graphs/}}

% ── Geometry ────────────────────────────────────────────────────
\geometry{
	a4paper,
	top=2.5cm, bottom=2.5cm,
	left=3cm, right=3cm
}

% ── Colors ─────────────────────────────────────────────────────
\definecolor{fxnavy}{RGB}{0,51,102}
\definecolor{fxgray}{RGB}{90,90,90}
\definecolor{fxlight}{RGB}{230,240,255}
\definecolor{warnyellow}{RGB}{255,200,0}
\definecolor{okgreen}{RGB}{0,128,64}

% ── Cover page colors ───────────────────────────────
\definecolor{docblue}{RGB}{22,74,140}
\definecolor{doclight1}{RGB}{213,227,245}
\definecolor{doclight2}{RGB}{231,240,251}
\definecolor{docgrey}{RGB}{90,90,90}
\usepackage[colorlinks = true, urlcolor = black, linkcolor = red]{hyperref}
%\usepackage{tikz} % For the frame
%\usetikzlibrary{calc}

% ── Hyperlinks ───────────────────────────────────────────────────
\hypersetup{
	colorlinks=true,
	linkcolor=fxnavy,
	citecolor=fxnavy,
	urlcolor=fxnavy,
	pdftitle={FOREX Volatility Regime Model},
	pdfauthor={Armand ZEWU}
}

% ── Header / Footer ───────────────────────────────────────
\pagestyle{fancy}
\fancyhf{}
\fancyhead[L]{\small\textcolor{fxgray}{FOREX Volatility Regime Model}}
\fancyhead[R]{\small\textcolor{fxgray}{\thepage}}
\renewcommand{\headrulewidth}{0.4pt}

% ── Spacing ───────────────────────────────────────────────────
\onehalfspacing

% ── Boxes ─────────────────────────────────────────────────────
\tcbuselibrary{skins,breakable}
\newtcolorbox{fxbox}[1][]{
	enhanced, breakable,
	colback=fxlight, colframe=fxnavy,
	fonttitle=\bfseries\color{fxnavy},
	title={#1}, arc=3pt
}
\newtcolorbox{limitbox}{
	enhanced, breakable,
	colback=yellow!10, colframe=orange!70!black,
	fonttitle=\bfseries, title={Documented limitation},
	arc=3pt
}
\newtcolorbox{corrbox}[1][Error identified and corrected]{
	enhanced, breakable,
	colback=red!5, colframe=red!60!black,
	fonttitle=\bfseries, title={#1},
	arc=3pt
}

\newtcolorbox{econbox}[1][Economic interpretation]{
	enhanced, breakable,
	colback=okgreen!8, colframe=okgreen,
	fonttitle=\bfseries, title={#1},
	arc=3pt
}



% ── Notation ─────────────────────────────────────────────────────
\newcommand{\R}{\mathbb{R}}
\newcommand{\Prob}{\mathbb{P}}
\newcommand{\E}{\mathbb{E}}

\setlength{\headheight}{13.59999pt}

% ---- Title page content ("SKDE-notat" style) ----
\newcommand{\titrerapport}{Quantitative modeling \\of EUR/USD exchange rate volatility}
%\newcommand{\soustitrerapport}{Quantitative modeling \\of EUR/USD exchange rate volatility}}
\newcommand{\resumerapport}{A quantitative model exploring the Forex market (EUR/USD),			capable of measuring an indicator \\of \guillemotleft~fear~\guillemotright\ /			volatility that explains returns and detects regimes	(low / normal / high).}

% ═════════════════════════════════════════════════════════════════
\begin{document} 
% ═════════════════════════════════════════════════════════════════

% ── Cover page (integrated FX-Cover code, A4 scale 21x29.7) ──
\newgeometry{margin=0cm}
\pagestyle{empty}
\thispagestyle{empty}
\noindent


\begin{titlepage}
	
	
	\begin{tikzpicture}[remember picture, overlay]
		
		%\IfFontExistsTF{Montserrat}{\fontspec{Montserrat}}{}
		
		% ---- Decorative background ----
		\fill[doclight1] ($(current page.north east)+(-3.0,1.4)$) circle (5.8cm);
		\fill[doclight2] ($(current page.south west)+(3.4,-2.0)$) circle (6.4cm);
		\fill[doclight1] ($(current page.south east)+(-0.6,2.6)$) circle (3.6cm);
		
		% ---- Title ----
		\node[anchor=north west, text width=15cm, align=left,
		text=docblue, font=\fontsize{27}{32}\selectfont\bfseries]
		at ($(current page.north west)+(2.2,-4.3)$) {\titrerapport};
		
		% ---- Subtitle ----
		%	\node[anchor=north west, text width=15cm, align=left,
		%	text=black, font=\Large]
		%	at ($(current page.north west)+(2.2,-8.9)$) {\soustitrerapport};
		
		% ---- Abstract ----
		\node[anchor=north west, text width=10cm, align=justify,
		font=\normalsize] (resume)
		at ($(current page.north west)+(2.2,-10.0)$) {\resumerapport};
		
		% ---- Rule + note under the abstract ----
		\draw[docgrey, line width=0.4pt]
		($(resume.south west)+(0,-0.35)$) -- ($(resume.south east)+(0,-0.35)$);
		\node[anchor=north east, font=\small]
		at ($(resume.south east)+(0,-0.55)$) {Quantitative research report --- 2026};
		
		% ---- Author block (vertical rule on the left) ----
		\draw[docblue, line width=1pt]
		($(current page.south east)+(-8.6,4.9)$) -- ($(current page.south east)+(-8.6,7.4)$);
		\node[anchor=north west, text width=7cm, align=left, font=\normalsize]
		at ($(current page.south east)+(-8.3,7.4)$)
		{ZEWU Yawo Armand Isaac\\
			\textit{Undergraduate Student in International Economics}\\
		};
		
		% ---- Footer: logo / institution ----
		%\foreach \x/\y/\c in {0/0.13/docblue, 0.32/0.13/docblue!55,
			%	0/-0.15/docblue!55, 0.32/-0.15/docblue!25}{
			%	\fill[\c] ($(current page.south west)+(2.2+\x,1.95+\y)$) circle (0.045);
			%}
		
		\node[anchor=west, text=docblue, font=\bfseries\normalsize]
		at ($(current page.south west)+(2.2,1.95)$) {Faculty of Economics and Management Sciences --- University of Lomé};
		\node[anchor=north west, text=docblue, font=\small]
		at ($(current page.south west)+(2.2,1.55)$)
		{Department of Economics};
		
	\end{tikzpicture}
	
\end{titlepage}	

\clearpage
\restoregeometry
\pagestyle{fancy}

% ── Abstract ──────────────────────────────────────────────────────
\begin{titlepage}
	\begin{center}
		
		\begin{mdframed}[backgroundcolor=fxlight,linecolor=fxnavy,
			linewidth=1pt,innerleftmargin=12pt,
			innerrightmargin=12pt,innertopmargin=10pt,
			innerbottommargin=10pt]
			\textbf{Abstract.}
			This document brings together the entirety of a quantitative
			modeling project on EUR/USD exchange rate volatility, conducted in
			two stages: (a) measuring, in a statistically testable way, a
			shift in market volatility, rather than describing it
			qualitatively; (b) deriving an operational use from it --- a
			price range adjusted to the market regime and real-time
			regime detection. A GARCH(1,1) with Student's~$t$
			innovations ($\hat\nu \approx 7{,}04$) is specified and validated on
			daily log returns (FRED, \texttt{DEXUSEU}, 2010--2026, $n=4\,175$).
			The model's temporal stability is tested over five equal
			sub-periods: point estimates vary, but a
			likelihood-ratio test does not reject the hypothesis of a
			single set of parameters ($p=0{,}277$). A baseline regression
			$r(t+1) = \alpha + \beta_1 \Delta Y(t) + \beta_2 r(t) + u(t)$
			detects no predictive power on the \emph{direction} of the
			return, a result robust to the addition of control variables
			(rate differential, DXY proxy). By contrast, the level of
			conditional volatility $\sigma(t)$ explains a
			significant share of the \emph{magnitude} of the future return
			($R^2 = 0{,}090$, $p<0{,}001$) --- a result of internal consistency
			of the GARCH, not a new discovery. A three-state Gaussian HMM
			and then a Markov-Switching GARCH ($K=2$, specification of
			Haas, Mittnik \& Paolella, 2004a) yield persistent regime
			detection (average durations of 72 to 404 days depending on the
			regime and the model). The resulting operational price range
			is validated by an expanding-window backtest of 3,424
			one-day-ahead forecasts over 2013--2026: empirical coverage
			of 95.50\% against a nominal target of 95\% (one-sided
			binomial test, $p=0{,}919$) --- the project's first validation
			with genuine statistical power. A dedicated out-of-sample
			robustness test for the COVID-19 crisis (January--April 2020,
			model re-estimated on strictly prior data) complements this
			result without contradicting it. Each step is documented with its
			errors, corrections, and acknowledged limitations, in the
			interest of full methodological transparency.
		\end{mdframed}
		
	\end{center}
\end{titlepage}

% ── Table of contents ───────────────────────────────────────────
\tableofcontents
\newpage

\begin{center}
	{\LARGE\bfseries Author's Note}
\end{center}
\vspace{0.8cm}

I first encountered risk aversion in microeconomics --- a stable individual preference within the framework of utility theory. But financial markets generally do not behave this way, as if risk perception were constant: volatility rises sharply, investor sentiment shifts, and markets alternate between periods of calm and panic.

This led me to ask a question that microeconomics alone cannot answer: \textit{\textbf{can this perception of risk, variable over time and at the scale of the market — which is closer to behavioral finance than to classical risk aversion — be measured quantitatively on the foreign exchange (FX) market, the most liquid market in the world? And, if so, can it help explain returns and produce a usable price range, rather than a purely qualitative analysis?}}

This project is also an opportunity to apply the skills acquired in Econometrics and Decision Statistics studied during my Bachelor's degree in International Economics at the FASEG of the University of Lomé.
It naturally connects with my courses in International Trade and International Monetary Economics, since financial markets sit at the intersection of these two fields.

I want to stress that the primary aim of this modeling exercise was not to build a Forex trading tool, but the results obtained, particularly the results on the price range, can be used for an analytical strategy on these markets.

\newpage
% ═════════════════════════════════════════════════════════════════
\section{Introduction}
\label{sec:intro}
% ═════════════════════════════════════════════════════════════════

\subsection{Project objective}

This project builds a quantitative model for the foreign exchange market
(EUR/USD as a case study), with a dual ambition:

\begin{itemize}
	\item[\textbf{(a)}] To measure, in a quantitative and statistically
	testable way, a behavioral shift in the market
	(fear~/~risk aversion), rather than describing it
	qualitatively as part of the behavioral finance
	literature does.
	\item[\textbf{(b)}] To derive a potential operational use from it: a
	decision-support signal --- not a certain prediction --- and
	real-time market regime detection (low/normal/high volatility (stress-panic)).
\end{itemize}





\subsection{Theoretical starting point}

\textbf{\textit{	The project's founding question is: does the future price $P(t+1)$ depend on the present price $P(t)$?}}

\begin{itemize}
	\item On a market as liquid as EUR/USD spot, the
	\textbf{direction} of the price approximately follows a random
	walk (weak-form efficient market hypothesis) --- nearly
	unpredictable from price history alone.
	\item By contrast, the \textbf{magnitude~/~volatility} of the
	movement is partially predictable, via the volatility
	\emph{clustering} phenomenon: a turbulent period tends to be
	followed by a turbulent period.
	\item Operational consequence: aim for a \textbf{price range} (a confidence interval based on volatility), not a fixed price.
\end{itemize}

These two results will be established empirically and separately in this document: the absence of directional predictability
(section~\ref{sec:baseline}, confirmed in section~\ref{sec:controls}) and the predictability of magnitude (section~\ref{sec:magnitude}).

\subsection{Model specification --- iterations}
\label{sec:spec-iterations}


The initial proposed equation was
$P(t+1) = X(t) + Y(t) + u(t)$, 

where:

\begin{itemize}
	\item $X$ is the price at time $t$,
	\item	$Y$ a market fear sentiment index, 
	\item and $u$ an error term. 
\end{itemize}

\newpage

Three successive corrections were needed before arriving at an estimable specification:

\begin{enumerate}
	\item \textbf{Dimensional inconsistency.} It is impossible
	to add a price and a fear index without a weighting
	coefficient: the correct form requires a $\beta$,
	$P(t+1) = P(t) + \beta \cdot Y(t) + u(t)$.
	\item \textbf{Non-stationarity of the price level.} Exchange
	rates are close to a unit-root process. Modeling the level exposes one to a spurious regression
	\citep{grangernewbold1974}. The model therefore operates on the
	\textbf{logarithmic return}, $r(t+1) = \ln\!\big(P(t+1)/P(t)\big)$.
	\item \textbf{Simultaneity~/~causality.} A fear index computed from contemporaneous option prices could be a consequence of the price, not its cause. The explanatory variable must therefore be lagged in time, or its causality tested
	explicitly.
\end{enumerate}

\vspace{0.8cm}

The working specification adopted is:

\begin{center}
	\begin{keyeq}
		\begin{equation}
			r(t+1) = \alpha + \beta_1 \, \Delta Y(t) + \beta_2 \, r(t)
			+ \sum_i \gamma_i \, \text{Control}_i(t) + u(t),
			\label{eq:master-spec}
		\end{equation}
	\end{keyeq}
\end{center}



\vspace{0.4cm}

where:

\begin{itemize}
	\item $\Delta Y(t)$ is the change in the fear~/~volatility variable. It is defined concretely in section~\ref{sec:garch}: it is the standardized residual of a GARCH(1,1) with Student's~$t$ innovations, calibrated on EUR/USD log returns (not its level), 
	\item $r(t)$ captures any residual momentum or \emph{mean-reversion}, and the controls include the interest rate differential and a dollar proxy (DXY).
\end{itemize} 


\subsection{Choice of the fear~/~volatility variable}

The usual reference indices were ruled out: the CVIX (Deutsche
Bank) and the JPMorgan VXY Global are proprietary products beyond
reach; the equity VIX (S\&P~500), considered as a proxy, was rejected because it captures a risk specific to U.S. equities,
potentially disconnected from FX --- the VIX
$\rightarrow$ FX transmission chain is indirect and unstable over time, and
using it would break the coherence of the object of study (measuring a proxy from another market rather than FX itself).

The solution adopted is an \textbf{internally constructed FX conditional
	volatility}, via a GARCH(1,1) calibrated directly on EUR/USD log returns --- not a simple rolling standard-deviation window.
General principle: an imperfect but well-targeted measure (the
volatility of the right market) outperforms a sophisticated but poorly
targeted measure (the volatility of another market).

\subsection{Data sources}
\label{sec:data-sources}

\begin{table}[H]
	\centering
	\caption{Data sources used}
	\label{tab:data-sources}
	\begin{tabular}{lll}
		\toprule
		Series & Source & Role \\
		\midrule
		EUR/USD (\texttt{DEXUSEU}) & FRED & Price $P(t)$, return $r(t)$ \\
		Fed Funds rate (\texttt{DFF}) & FRED & Control --- rate differential \\
		ECB deposit rate (\texttt{ECBDFR}) & FRED & Control --- rate differential \\
		Broad dollar index (\texttt{DTWEXBGS}) & FRED & Control --- DXY proxy \\
		EUR/USD OHLC (external reference) & Pound Sterling Live & Independent verification \\
		\bottomrule
	\end{tabular}
\end{table}

\vspace{0.4cm}

The choice of FRED as the primary source follows the failure of
\texttt{yfinance} on the \texttt{EURUSD=X} ticker
(section~\ref{sec:garch-data}). A convention point deserves to be
fixed from the introduction: \texttt{DEXUSEU} is quoted in dollars per
euro (standard EUR/USD market convention, the euro being the base currency) ---  
A second clarification, reused several times in this document:
\texttt{DEXUSEU} is the \emph{noon
	buying rate} published by the Federal Reserve (H.10 release) in New
York, not a market closing price --- a distinction that explains
part of the discrepancies observed later (section~\ref{sec:covid}) between
this project and independent OHLC sources.

\subsection{Document roadmap}

Table~\ref{tab:roadmap} summarizes the structure of the document and,
implicitly, the logical order in which the questions were
raised and closed over the course of the project.

\begin{table}[H]
	\centering
	\small
	\caption{Roadmap}
	\label{tab:roadmap}
	\begin{tabular}{@{}p{0.9cm}p{6.3cm}p{7.2cm}@{}}
		\toprule
		\S{} & Step & Question asked \\
		\midrule
		\ref{sec:garch} & GARCH(1,1) & How to cleanly isolate a volatility shock
		from mere inertia? \\
		\ref{sec:garch-stability} & Temporal stability of the GARCH & Is the
		single model a stable description, or an average that
		masks regimes? \\
		\ref{sec:baseline} & Baseline regression & Does $\Delta Y(t)$
		predict the \emph{direction} of $r(t+1)$? \\
		\ref{sec:controls} & Regression with controls & Does this null result
		survive the addition of the rate differential and the DXY? \\
		\ref{sec:magnitude} & Magnitude regression & Do $\Delta Y(t)$ or
		$\sigma(t)$ predict the \emph{size} of $r(t+1)$? \\
		\ref{sec:regime-hmm} & Regime detection (HMM) & Can days be
		classified into volatility regimes, without re-estimating the
		GARCH by regime? \\
		\ref{sec:msgarch} & Markov-Switching GARCH & What changes when
		$\omega,\alpha,\beta$ are re-estimated separately by regime? \\
		\ref{sec:price-range} & Operational price range & What is
		objective (b) of the project actually worth, concretely? \\
		\ref{sec:rolling-test} & Rolling chain test & How does the
		range behave day after day, under real
		conditions? \\
		\ref{sec:covid} & Out-of-sample robustness (COVID) & What
		happens during a major shock, with a model that has
		never seen it? \\
		\ref{sec:backtest} & Full backtest (2013--2026) & Is the model
		calibrated, with genuine statistical power? \\
		\ref{sec:conclusion} & Conclusion & What does this project establish,
		and what does it leave open? \\
		\bottomrule
	\end{tabular}
\end{table}

%	\textbf{Note on the order chosen.} The section on the temporal
%	stability of the GARCH (\ref{sec:garch-stability}) is placed
%	immediately after the GARCH specification, even before the
%	baseline regression: this is the logical order in which this
%	question was raised --- \emph{before} the choice of a regime
%	detection method (HMM vs Markov-Switching GARCH) --- even though,
%	in the project's methodology log, its written confirmation came
%	later. The placement of the figures follows exactly
%	the arbitration already set in the dedicated working document
%	(\emph{Graphs\_Ways.md}) and is not modified here.

% ═════════════════════════════════════════════════════════════════
\newpage
\section{The GARCH(1,1) --- Specification and Validation}
\label{sec:garch}
% ═════════════════════════════════════════════════════════════════

\subsection{Objective}

Specify and validate a GARCH(1,1) on the daily log returns
of EUR/USD, in order to
\begin{itemize}
	\item (a) isolate the volatility shock from its mechanical
	inertia (\emph{clustering}), 
	\item (b) obtain a series of standardized
	residuals usable as a "fear shock" input variable
	for the regression on the return (section~\ref{sec:baseline})
\end{itemize}
.

\subsection{Data pipeline}
\label{sec:garch-data}



The loader (\texttt{data\_loader\_fred.py}) retrieves the
\texttt{DEXUSEU} series directly from FRED's CSV
endpoint: \textbf{4,176 daily observations} successfully obtained (2010-01-01 to 2026-09-11).


%\textbf{Zero-return check.} 39 of the 4,175 daily log
%returns are exactly zero ($\approx 0{,}93\,\%$ of
%observations). A dispersion test around known bank
%holidays (which would suggest a stale value repeated by
%FRED) shows no visible calendar pattern: the dates are
%scattered. Conclusion: consistent with FRED's rounding to 4 decimals
%on low-volatility days, not a data quality issue ---
%the series is left intact.

\subsection{Diagnostics of the return series (before GARCH)}

Log returns (scaled $\times 100$), full sample,
$n = 4\,175$:

\begin{table}[H]
	\centering
	\caption{Descriptive statistics of EUR/USD log returns}
	\begin{tabular}{lr}
		\toprule
		Statistic & Value \\
		\midrule
		Mean & $-0{,}0052\,\%$ \\
		Standard deviation & $0{,}518\,\%$ \\
		Minimum & $-2{,}672\,\%$ \\
		Maximum & $3{,}064\,\%$ \\
		\bottomrule
	\end{tabular}
\end{table}

A mean close to zero is consistent with weak-form market
efficiency (no exploitable mean drift). The minimum and
maximum, well beyond several standard deviations from the mean,
provide an initial visual indication of fat tails,
formally confirmed later (\S\ref{sec:garch-fattails}).

\vspace{0.4cm}

\begin{figure}[H]
	\centering
	\includegraphics[width=1\linewidth]{01_returns_series.png}
	\caption{Daily log return series, EUR/USD, 2010--2026.}
	\label{fig:returns-series}
\end{figure}

\newpage

\subsection{GARCH(1,1) --- Normal distribution (first specification)}

\begin{center}
	\begin{keyeq}
		
		\begin{equation}
			\sigma^2(t) = \omega + \alpha \, \varepsilon^2(t-1) + \beta \, \sigma^2(t-1)
		\end{equation}
	\end{keyeq}
\end{center}


\vspace{0.8cm}

\begin{table}[H]
	\centering
	\caption{GARCH(1,1), Normal innovations}
	\begin{tabular}{lrr}
		\toprule
		Parameter & Estimate & Significance \\
		\midrule
		$\mu$ & $-0{,}00588$ & $p=0{,}394$, not significant \\
		$\omega$ & $0{,}001446$ & $p=0{,}031$ \\
		$\alpha_1$ & $0{,}0350$ & $p \approx 1{,}6\times10^{-9}$ \\
		$\beta_1$ & $0{,}9595$ & $p<0{,}001$ \\
		\midrule
		\multicolumn{3}{l} {Log-likelihood $=-2900{,}17$ ; AIC $=5808{,}34$ ; BIC $=5833{,}69$} \\
		\bottomrule
	\end{tabular}
\end{table}


\begin{itemize}
	\item $\mu$ not significant is consistent with the random-walk
	hypothesis for the direction of returns. 
	\item $\alpha+\beta = 0{,}9945 < 1$ (stationarity condition satisfied).
	\item Implied shock half-life: $\ln(0{,}5)/\ln(0{,}9945) \approx 126$ trading days ($\approx$~6 months).
\end{itemize}


\newpage

\subsection{Diagnostics --- Normal model}

\begin{corrbox}
	The initial diagnostic script called
	\texttt{result.arch\_lm\_test(lags=10)} \textbf{without}
	\texttt{standardized=True}. This test then runs on the
	raw residuals of the mean model, not on the residuals standardized by the
	GARCH --- which almost always rejects strongly for a series of
	financial returns, not because the GARCH is misspecified, but
	because this is expected for any return series before
	volatility modeling. Seemingly contradictory result:
	\begin{center}
		\begin{tabular}{| p{4.7cm} | p{4.7cm} |}
			\hline
			Test & Result \\
			\hline
			Ljung-Box on squared standardized residuals (lags 10, 20) & $p=0{,}933$ / $0{,}591$ --- no residual ARCH effect \\
			\hline
			ARCH-LM (raw residuals, bug) & $p \approx 0{,}0000$ --- apparent rejection \\
			\hline
		\end{tabular}
	\end{center}
	\textbf{Correction:} re-run with \texttt{standardized=True}.
	Result: statistic $=4{,}14$, $p=0{,}941$ --- consistent with the
	Ljung-Box test. The two tests agree once correctly run
	on the standardized series. No actual misspecification of the
	model: the contradiction was a test configuration error,
	not a flaw in the model.
\end{corrbox}

\newpage 

\subsubsection{Fat tails}
\label{sec:garch-fattails}

Excess kurtosis of standardized residuals: \textbf{1.620} (Normal $=0$).
The QQ-plot against the Normal distribution shows a clear
deviation in the tails.
Range of standardized residuals: $-5{,}91$ to $+5{,}02$ --- under a
true Normal distribution, values at $\pm5$ standard deviations would
be nearly impossible at this sample size.

\vspace{0.8cm}

\begin{figure}[H]
	\centering
	\begin{subfigure}{0.45\linewidth}
		\includegraphics[width=1\linewidth]{02_returns_histogram_vs_normal.png}
		\caption{Histogram of returns against the Normal distribution}
	\end{subfigure}
	\hfill
	\begin{subfigure}{0.45\linewidth}
		\includegraphics[width=1\linewidth]{05_qqplot_normal_model.png}
		\caption{QQ-plot, Normal model}
	\end{subfigure}
	\caption{Evidence of fat tails under the Normal specification.}
	\label{fig:fattails-normal}
\end{figure}

\textbf{Decision:} re-estimate with a Student's~$t$
innovation distribution.

\vspace{0.6cm}

\subsection{GARCH(1,1) --- Student's $t$ distribution (second specification)}

\begin{table}[H]
	\centering
	\caption{GARCH(1,1), Student's $t$ innovations}
	\begin{tabular}{lrr}
		\toprule
		Parameter & Estimate & Significance \\
		\midrule
		$\mu$ & $-0{,}00578$ & $p=0{,}368$, not significant \\
		$\omega$ & $0{,}000651$ & $p=0{,}099$, \textbf{not significant} (vs $p=0{,}031$ under Normal) \\
		$\alpha_1$ & $0{,}0374$ & $p \approx 3{,}3\times10^{-12}$ \\
		$\beta_1$ & $0{,}9611$ & $p<0{,}001$ \\
		$\nu$ & $7{,}0414$ & $p \approx 4{,}1\times10^{-22}$ \\
		\midrule
		\multicolumn{3}{l}{Log-likelihood $=-2818{,}39$ ; AIC $=5646{,}78$ ; BIC $=5678{,}46$} \\
		\bottomrule
	\end{tabular}
\end{table}

\begin{itemize}
	\item $\alpha+\beta = 0{,}9985$ (vs $0{,}9945$ under Normal) $\Rightarrow$
	\item Implied half-life $\approx 462$ days ($\approx$~18 months), nearly 4 times longer than under the Normal specification.
\end{itemize}

\vspace{0.3cm}

\begin{limitbox}
	This genuine sensitivity of persistence $\alpha+\beta$ to
	the distributional assumption (0.9945 vs 0.9985) is not a
	rounding artifact --- it is treated as a range, not a single
	figure, until the sub-period stability test
	(section~\ref{sec:garch-stability}). $\omega$ loses its significance
	under this specification --- flagged, not dismissed.
\end{limitbox}

\vspace{0.6cm}

\subsection{Model comparison and selection}


\begin{table}[H]
	\centering
	\caption{Normal vs Student's $t$ comparison}
	\begin{tabular}{lrr}
		\toprule
		Criterion & Normal & Student's $t$ \\
		\midrule
		AIC & 5808.34 & \textbf{5646.78} \\
		BIC & 5833.69 & \textbf{5678.46} \\
		Log-likelihood & $-2900{,}17$ & \textbf{$-2818{,}39$} \\
		Ljung-Box $p$ (lag 10) & 0.933 & 0.972 \\
		ARCH-LM $p$ (standardized) & 0.941 & 0.976 \\
		Excess kurtosis (empirical) & 1.620 & 1.831 \\
		\bottomrule
	\end{tabular}
\end{table}

AIC and BIC clearly favor the Student's~$t$ specification
($\Delta$AIC $\approx 161{,}6$, $\Delta$BIC $\approx 155{,}2$ ---
both well beyond conventional model comparison
thresholds).

\vspace{0.5cm}

\begin{corrbox}
	An initial reading of the Student's~$t$ diagnostics flagged the
	empirical excess kurtosis of 1.831 as "still too high, tails
	not absorbed" --- but this comparison used 0 as the
	reference, a benchmark valid only for a
	\textbf{Normal} model. For a Student's~$t$ model with $\nu$ degrees of
	freedom, the theoretical excess kurtosis equals
	$6/(\nu-4)$. With $\nu = 7{,}0414$:
	$6/3{,}0414 \approx \mathbf{1{,}973}$.
	The empirical value (1.831) is close to this theoretical value
	(1.973) --- this is a \textbf{good} result, indicating that the estimated
	$\nu$ correctly describes the actual shape of the residual
	distribution. The initial flag was a diagnostic threshold
	error (comparison against the wrong reference), not an actual
	flaw in the model.
\end{corrbox}

\newpage

\begin{figure}[H]
	\centering
	\includegraphics[width=1\linewidth]{06_qqplot_studentt_model.png}
	\caption{QQ-plot, Student's $t$ model --- visual confirmation of the correction above.}
	\label{fig:qqplot-studentt}
\end{figure}

\subsection{Final decision}

\textbf{Model selected: GARCH(1,1) with Student's~$t$ innovations
	($\nu \approx 7{,}04$).} The Ljung-Box and standardized
ARCH-LM tests both confirm the absence of residual conditional
heteroskedasticity. The empirical kurtosis closely matches
the model-implied kurtosis. AIC and BIC both favor this
specification over the Normal alternative, decisively.

The series of standardized residuals from this model constitutes the project's
operational $\Delta Y(t)$ variable, used as the volatility shock
input for the rest of the document.

\newpage

\subsection{Economic interpretation}

\begin{econbox}[What $\omega$, $\alpha$, $\beta$ mean]
	
	The chosen model says: tomorrow's expected volatility is a
	weighted combination of 
	
	\begin{itemize}
		\item a long-run reference level,
		\item the size of today's price shock
		\item the volatility already observed today
	\end{itemize}
	
	\textbf{$\alpha \approx 0{,}037$}:
	only 3.7\,\% of a single day's shock feeds directly through
	into tomorrow's volatility estimate --- an isolated piece of bad news
	does not, by itself, strongly move the market's anticipated
	volatility. \textbf{$\beta \approx 0{,}961$}: about 96\,\% \\of
	today's volatility level persists into tomorrow --- by far the
	dominant term. \textbf{Volatility is mostly inertia}, not a
	reaction to an isolated event.
\end{econbox}

\vspace{2cm}

\begin{econbox}[What persistence and half-life mean]
	
	$\alpha+\beta \approx 0{,}9985$, i.e. a half-life of about 462
	trading days --- nearly \textbf{18 months} for a volatility shock
	to lose half its intensity. In plain terms: when EUR/USD
	volatility spikes (after a surprise central bank decision or
	a geopolitical shock), the elevated uncertainty does \textbf{not}
	dissipate within a few days or weeks --- it fades slowly, over more
	than a year, before returning close to its reference level.
	\textbf{Important caveat}: section~\ref{sec:garch-stability}
	shows that this 18-month figure is itself unstable depending on the
	sub-period (from 2--3 months to about a year). The honest statement is:
	EUR/USD volatility shocks take between 2 months and more than
	a year to fade, depending on the period considered --- not a single
	figure to cite.
\end{econbox}

\newpage

\begin{econbox} [What fat tails ($\nu \approx 7$) mean --- a concrete illustration]
	
	The most extreme daily move in the sample
	($\approx -2{,}67\,\%$) corresponds to a standardized shock
	of about 5.9 standard deviations under the fitted model.
	\begin{itemize}
		\item \textbf{Under the Normal model}: probability of about
		$1{,}7\times 10^{-9}$ --- an event with \textbf{1 chance in
			576 million trading days}, i.e. once every
		\textbf{2.3 million years}. Since the project's sample only
		covers 16 years, nonetheless observing this move shows that
		the Normal model tells an implausible story about
		the market's actual behavior.
		\item \textbf{Under the chosen Student's~$t$ model} ($\nu \approx 7{,}04$):
		probability of about $0{,}00029$ --- a frequency of about
		\textbf{once every 14 years}. A plausible frequency for a
		"large but not unprecedented" shock over a 16-year sample
		that includes central bank surprises and acute
		risk-aversion episodes.
	\end{itemize}
	In plain terms: a model assuming Normal shocks would lead one to treat
	large EUR/USD moves as near-impossible "black swan" events
	--- only to be caught off guard when they occur, a well-documented
	failure mode in risk management. The Student's~$t$ model tells a
	much more realistic story: large moves are rare, but not so rare
	that a market participant could assume they will
	never happen.
\end{econbox}

\begin{figure}[H]
	\centering
	\includegraphics[width=1\linewidth]{09_tail_probability_comparison.png}
	\caption{Probability of the sample's most extreme move, under the two specifications --- the gap between "once every 2.3 million years" and "once every 14 years."}
	\label{fig:tail-probability}
\end{figure}


\vspace{0.6cm}

\begin{econbox}[For a market observer]
	
	EUR/USD volatility moves in slow, persistent waves rather
	than by daily chance; it takes months
	(with genuine uncertainty about the exact number) to calm down after
	a shock; and it produces large, sudden moves
	noticeably more often than a "textbook Normal" view of the
	market would suggest. In practice, this justifies treating
	current volatility as an informative, slowly evolving signal
	--- and directly motivates the project's next step: building a
	volatility-sensitive price range
	(section~\ref{sec:magnitude}, then section~\ref{sec:price-range})
	rather than relying on a single fixed estimate of the
	possible size of a move.
\end{econbox}

\newpage

\subsection{Known limitations of this step}

\begin{limitbox}
	\begin{itemize}
		\item $\alpha+\beta$ is strongly sensitive to the
		distributional assumption (0.9945 vs 0.9985) --- the implied
		persistence and half-life must be reported as a range,
		not a single confident figure, pending a stability test by
		sub-periods (section~\ref{sec:garch-stability}).
		\item $\omega$ is not statistically significant under the
		Student's~$t$ specification --- flagged, not resolved at this stage.
		\item This entire specification relies solely on
		daily noon-rate data (\texttt{DEXUSEU} is the Federal
		Reserve's noon reference rate, not a market
		closing price --- \S\ref{sec:garch-data}), with no
		intraday realized volatility --- a data-access limitation
		documented again later in the project
		(section~\ref{sec:covid}).
	\end{itemize}
\end{limitbox}

\newpage

% ═════════════════════════════════════════════════════════════════
\section{Temporal stability of the GARCH}
\label{sec:garch-stability}
% ═════════════════════════════════════════════════════════════════

\subsection{Why this test precedes the choice of regime \\ detection method}

This test validates the \textbf{measurement tool} itself (the GARCH model),
not its operational use (price ranges, regime detection).

It asks whether the single GARCH(1,1) calibrated on the full 2010--2026 sample
--- the one used throughout the document --- is a
stable description of volatility dynamics over time, or an average that poorly represents each sub-period taken in isolation.

This test was deliberately conducted \textbf{before} choosing a regime detection method (HMM vs Markov-Switching GARCH), since the stability result was itself meant to inform this choice: applying a simple HMM to the \textbf{$\sigma(t)$} of an unstable global GARCH would inherit that instability as a hidden bias, whereas a genuinely unstable GARCH is itself an argument for a Markov-Switching GARCH (which re-estimates $\alpha$, $\beta$, $\omega$ by regime) rather than for a
simpler two-step approach.

\subsection{Method}

The full sample has 4,175 observations. 5 periods divide it exactly: $4175/5 = 835$. This choice is motivated solely by this arithmetic property.

The split is \textbf{chronological, non-overlapping, of equal
	size} --- not defined by economically motivated breaks
("before/after COVID", "before/after the euro-area debt
crisis"). A deliberate choice to avoid selection bias:
choosing break dates based on known economic events would risk
unknowingly selecting dates that confirm an
already-assumed narrative rather than neutrally testing stability.

\vspace{0.3cm}

\begin{table}[H]
	\centering
	\caption{Periods used}
	\begin{tabular}{|p{2.7cm}|c|p{3.2cm}|}
		\hline
		Period & Date range & $n$ \\
		\hline
		1 & 2010-01-05 to 2013-05-01 & 835 \\
		\hline
		2 & 2013-05-02 to 2016-08-29 & 835 \\
		\hline
		3 & 2016-08-30 to 2020-01-07 & 835 \\
		\hline
		4 & 2020-01-08 to 2023-05-10 & 835 \\
		\hline
		5 & 2023-05-11 to 2026-09-11 & 835 \\
		\hline
	\end{tabular}
\end{table}

\vspace{0.3cm}

Each period's GARCH(1,1) with Student's~$t$ innovations is
estimated independently, using the specification already validated in
section~\ref{sec:garch}.

\subsection{Results}

\begin{table}[H]
	\centering
	\small
	\caption{Estimated parameters by period}
	\begin{tabular}{lrrrrrr}
		\toprule
		Period & $\omega$ & $\alpha$ & $\beta$ & $\nu$ & $\alpha+\beta$ & Half-life (days) \\
		\midrule
		1 & 0.004378 & 0.0283 & 0.9612 & 12.59 & 0.9895 & 65.5 \\
		2 & 0.000382 & 0.0386 & 0.9614 & 5.95 & \textbf{0.999991} & \textbf{80,366 (not significant --- \S\ref{sec:garch-stability-p2})} \\
		3 & 0.000233 & 0.0170 & 0.9809 & 10.58 & 0.9979 & 335.1 \\
		4 & 0.003341 & 0.0720 & 0.9192 & 6.89 & 0.9912 & 78.5 \\
		5 & 0.001878 & 0.0408 & 0.9494 & 5.54 & 0.9902 & 70.5 \\
		\bottomrule
	\end{tabular}
\end{table}

\vspace{0.6cm}

\begin{figure}[H]
	\centering
	\includegraphics[width=1\linewidth]{01_alphabeta_vs_halflife.png}
	\caption{$\alpha+\beta$ and implied half-life by period --- a direct visualization of the non-linearity that the table alone does not show.}
	\label{fig:alphabeta-halflife}
\end{figure}

\vspace{0.6cm}

\subsection{Diagnosis of the Period 2 anomaly}
\label{sec:garch-stability-p2}

Period 2's implied half-life ($\approx 80,366$ days,
$\approx 220$ years) is not economically meaningful and was
investigated before any interpretation.

\vspace{0.3cm}

\textbf{Root cause:} only $\omega$ is not statistically
distinct from zero in this sub-sample ($p=0{,}519$). The GARCH's
long-run (unconditional) variance equals $\omega/(1-\alpha-\beta)$.

When $\alpha+\beta$ is very close to 1 (here, 0.999991) and
$\omega$ is close to zero and imprecisely estimated, this ratio becomes
numerically unstable --- \textbf{\textit{a tiny change in the denominator produces
		a huge variation in the implied half-life}}. 

This is a numerical fragility
well known in GARCH models operating near the IGARCH boundary
($\alpha+\beta=1$ exactly), not a sign that the model
itself has failed.



\begin{figure}[H]
	\centering
	\includegraphics[width=0.9\linewidth]{05_calm_narrative_debunked.png}
	\caption{The "calm period" narrative for Period 2, compared against the stress-regime share observed in 2014--2015 by the MS-GARCH.}
	\label{fig:calm-narrative-debunked}
\end{figure}

\textbf{Conclusion on Period~2:} $\alpha$ and $\beta$ themselves
remain valid and interpretable (consistent with the other four
periods); the derived half-life, however, must be treated as not
meaningful and should never be cited alone without this explanation.


\subsection{Overall stability conclusion --- revised after a formal test}

Excluding the Period 2 half-life artifact (while keeping
its valid $\alpha, \beta$ estimates), the point
estimates vary substantially from one period to another ---
half-lives of about 65--80 days (Periods 1, 4, 5) to about
335 days (Period 3). This large gap in half-life corresponds to a
modest absolute gap in $\alpha+\beta$ itself (0.0067 to 0.0084) ---
the half-life formula, $\ln(0{,}5)/\ln(\alpha+\beta)$, is strongly
non-linear near the $\alpha+\beta=1$ boundary, so that
small differences in the raw parameter translate into large differences
in the derived half-life. Both facts are true at the same time and do not
cancel each other out: the underlying parameter gaps are modest in
absolute terms, and their practical consequence (how long a shock
takes to fade) is large.

\newpage

\begin{corrbox}
	The earlier version of our calculations led us to conclude, from the sole %(section~\ref{sec:garch-stability})
	table above, that "the single GARCH(1,1) on the full
	sample is not temporally stable." \textbf{This
		overstated what the data can actually support.} A
	likelihood-ratio test comparing the 5 independently fitted
	sub-period models (5 periods $\times$ 5
	parameters, including $\mu$ --- 25 parameters in total) to the single
	full-sample model (5 parameters) gives:
	\[
	LR = 2 \times \big[(-2806{,}77) - (-2818{,}39)\big] = 23{,}24,
	\qquad df = 20, \qquad p = 0{,}277.
	\]
	\[
	\text{AIC: } 5646{,}78 \text{ (single model)} \text{ vs } 5663{,}53 \text{ (5 separate periods)}
	\]
	\[
	\text{BIC: } 5678{,}46 \text{ (single model)} \text{ vs } 5821{,}95 \text{ (5 separate periods)}
	\]
	\textbf{At $p=0{,}277$, the null hypothesis of a single,
		stable set of parameters across the 5 periods cannot be rejected at any
		conventional threshold.} AIC and BIC both prefer the
	single model over the 5 separate models (lower is better for both
	criteria) --- the added flexibility of period-specific
	parameters does not offset its added complexity under either
	criterion.
\end{corrbox}

\vspace{0.6cm}

\begin{figure}[H]
	\centering
	\includegraphics[width=1\linewidth]{06_single_vs_5period_model.png}
	\caption{Single model vs 5 separate models --- the likelihood-ratio test ($LR=23{,}24$, $p=0{,}277$) as statistical evidence.}
	\label{fig:single-vs-5period}
\end{figure}

\newpage

\textbf{General conclusion:} the point estimates in the table
above do genuinely differ from one period to another, but
\textbf{this variation is not statistically established} as
genuine instability rather than sampling noise around a
single true set of parameters. The honest summary is: \textbf{\emph{point
		estimates vary from one sub-period to another; a formal
		likelihood-ratio test does not establish that this variation
		reflects genuine parameter instability rather than chance.}}
This is a markedly weaker claim than "the model is
unstable."

\subsection{Implication for the choice of regime detection method}

This result nonetheless informed the project's next step (regime
detection), but the reasoning must be stated cautiously: a
simple HMM applied to the $\sigma(t)$ generated by the single
full-sample GARCH \emph{could} inherit a distortion specific to
certain periods, but this document's formal test does not establish
that such a distortion is statistically real. The case for
a \textbf{Markov-Switching GARCH} therefore rests more solidly on
AIC (5633.27 vs 5646.78, favoring the MS-GARCH ---
section~\ref{sec:msgarch}) than on any "confirmed instability"
from this document --- and even there, \textbf{the BIC prefers the single-regime model} (5696.64 vs 5678.46). 

The honest framing is: the
MS-GARCH was pursued because it is the theoretically more
appropriate tool for regime-dependent dynamics and because it wins
on an information criterion, not because this stability test
proved the single-regime model to be inadequate.

\newpage

\subsection{Economic interpretation}

\begin{econbox}
	This document's revised conclusion has a direct economic
	reading, easy to miss amid the statistical detail: \textbf{there
		is no solid evidence that the persistence of EUR/USD volatility
		behaves differently across eras} --- a 2011 shock and a 2023
	shock fade at statistically indistinguishable rates, as far as
	this test can tell. For operational use, this is reassuring in a
	precise sense: the single full-sample GARCH is not clearly the
	wrong tool because "market behavior has fundamentally
	changed" --- the more mundane explanation (sampling
	variation in a genuinely noisy process) cannot be ruled out, and
	Occam's razor favors it given the formal test's result.
	
	At the same time, the corrected Period 2 narrative is a useful
	warning against another common error: reaching for a
	plausible economic story ("it was a calm period") to
	explain an inconvenient statistical result, without checking it against
	independent evidence. The story seemed right at the time; it
	turned out to be false once the project had better tools (the
	MS-GARCH's regime classification) to check it. The
	practical lesson: a numerically fragile estimate ($\omega$ close to
	zero) deserves a flagged \emph{statistical} explanation (the
	IGARCH-boundary fragility, which holds) rather than an improvised
	\emph{economic} explanation to make the figure less surprising
	(the calm-period narrative, which does not hold).
\end{econbox}

\vspace{0.5cm}
\begin{limitbox}
	\texttt{garch\_stability\_comparison.csv} keeps only the
	point estimates ($\omega, \alpha, \beta, \nu$) by period
	--- neither their standard errors nor their $p$-values. The $p$-values
	discussed above for Period 2 (e.g.\ $p=0{,}519$ for $\omega$)
	come from an interactive re-estimation of that period alone,
	run in the console rather than captured in the saved output of the
	script. A reader trying to \\reproduce this diagnostic from the
	CSV alone could not do so. This gap is flagged as such rather
	than corrected retroactively here.
\end{limitbox}

% ═════════════════════════════════════════════════════════════════
\newpage

\section{Baseline regression --- $r(t+1)$ on the GARCH volatility shock}
\label{sec:baseline}
% ═════════════════════════════════════════════════════════════════

\subsection{Objective}

Test whether the level shift in volatility estimated by the GARCH
explains the next day's EUR/USD return --- a \textbf{baseline}
specification, with no control variables, as decided in the
introduction (\S\ref{sec:spec-iterations}): the control variables
(rate differential, DXY) are deliberately left out at this stage, so as
not to confound the estimate of the volatility shock's own effect.

\subsection{Specification}


\begin{center}
	\begin{keyeq}
		\begin{equation}
			r(t+1) = \alpha + \beta_1 \, \Delta Y(t) + \beta_2 \, r(t) + u(t)
		\end{equation}
	\end{keyeq}
\end{center}		

\vspace{0.8cm}

where 
\begin{itemize}
	\item $r(t)$ is the daily EUR/USD log return ($\times 100$)
	\item $\Delta Y(t) = \ln\!\big(\sigma(t)/\sigma(t-1)\big)$ is the log change in the GARCH's conditional volatility.
\end{itemize}

\vspace{0.8cm} 

\textbf{Why $\sigma(t)$, not the standardized residual $z(t)$.} 

The GARCH produces two series with opposite statistical properties: 
\begin{itemize}
	\item conditional volatility $\sigma(t)$ is a persistent, slowly evolving \textbf{level} (autocorrelated by construction --- that is the whole point
	of the GARCH)~ 
	\item the standardized residual $z(t)$ is the cleaned \textbf{shock}, which behaves close to white noise if the GARCH is well specified
	(confirmed in section~\ref{sec:garch}: no residual ARCH effect).
\end{itemize}

\vspace{0.8cm} 

Taking
a "change" in $z(t)$ would amount to computing the change of an already-cleaned series --- not economically meaningful.
$\Delta Y(t)$, as a \emph{level-shift} signal,
requires the persistent series, hence $\sigma(t)$.

\newpage

\subsection{Data alignment}

$r(t+1)$ is constructed via a \emph{forward}
shift (\texttt{shift(-1)}) --- explicitly not a lag ---
documented in code comments to avoid reversing the
causal direction (the shock at $t$ explains the return at $t+1$, not the reverse).
Final sample after alignment and removal of edge
\texttt{NaN}s: \textbf{4,173 observations} (from the 4,175 original
log returns).

\subsection{Estimation method}

OLS via \texttt{statsmodels}, with \textbf{HAC
	(Newey-West) standard errors}, \texttt{maxlags=5}. This choice is motivated by the fact that
financial return residuals commonly exhibit
heteroskedasticity and/or residual autocorrelation, which would
otherwise understate standard errors and inflate apparent
significance --- a well-documented risk of false positives
in empirical finance.

\subsection{Results}

\begin{table}[H]
	\centering
	\caption{Baseline regression}
	\begin{tabular}{lrrrr}
		\toprule
		& Coef. & Std. err. & $z$ & $P>|z|$ \\
		\midrule
		const & $-0{,}0051$ & $0{,}008$ & $-0{,}648$ & $0{,}517$ \\
		$\Delta Y(t)$ & $-0{,}1546$ & $0{,}240$ & $-0{,}644$ & $0{,}519$ \\
		$r(t)$ & $0{,}0206$ & $0{,}017$ & $1{,}231$ & $0{,}218$ \\
		\midrule
		\multicolumn{5}{l}{$R^2 = 0{,}001$ ; adjusted $R^2 = 0{,}000$ ; $F = 0{,}976$ ($p=0{,}377$)} \\
		\bottomrule
	\end{tabular}
\end{table}

\textbf{Model not jointly significant} 
($p=0{,}377$).
\textbf{$\beta_1$ $(\Delta Y(t)$): $-0{,}1546$, $p=0{,}519$ --- not
	significant.} No predictive effect detected of the volatility
level shift on the direction of the next day's return.
\textbf{$\beta_2$ $(r(t))$: $0{,}0206$, $p=0{,}218$ --- not
	significant.} No momentum or mean reversion detected at a
one-day horizon. Residual diagnostics: Jarque-Bera
$p \approx 1{,}6\times 10^{-126}$ (strongly non-normal residuals,
excess kurtosis $\approx 4{,}8$) --- expected given the
fat-tailed nature of FX returns already documented
(section~\ref{sec:garch}).

\begin{figure}[H]
	\centering
	\begin{subfigure}{0.48\linewidth}
		\includegraphics[width=\linewidth]{01_scatter_rt1_vs_deltay.png}
		\caption{Scatter plot of $r(t+1)$ against $\Delta Y(t)$}
	\end{subfigure}
	\hfill
	\begin{subfigure}{0.48\linewidth}
		\includegraphics[width=\linewidth]{03_forest_plot_coefficients.png}
		\caption{\emph{Forest plot} of the coefficients}
	\end{subfigure}
	\caption{Visual absence of relationship, an immediate complement to the regression table.}
	\label{fig:baseline-results}
\end{figure}
\vspace{0.8cm}
\
\subsection{Interpretation}

This is a \textbf{null result on direction}, not a failed
model. It is consistent with the weak-form market efficiency
hypothesis stated in the introduction (\S\ref{sec:intro}): on a
market as liquid as EUR/USD spot, the \textbf{direction}
of the next day's return is expected to be nearly unpredictable from
the contemporaneous volatility level or the previous day's return alone.
This aligns with the well-documented empirical finding that
structural exchange rate models struggle to beat a random
walk at short horizons \citep{meeserogoff1983}.
\vspace{0.4cm}	
\textbf{What this does not call into question:} objective (a) of the project
--- quantitatively measuring a market fear~/~volatility shift
--- had already been achieved and validated in section~\ref{sec:garch},
independently of whether this shift predicts the direction
of the return. This regression tests one specific, falsifiable use
of this measure; that it returns a null result is itself a
valid, documented result, not a reason to add variables
until significance is achieved (a \emph{p-hacking} risk
explicitly avoided here).

\newpage
\subsection{What this means economically, even setting statistical significance aside}

\begin{econbox}
	Even \emph{if} these coefficients reflected a real relationship,
	would the effect be large enough to matter for someone
	trading or managing risk on EUR/USD? \textbf{\textit{The answer is no}}, on
	both counts.
	
	\vspace{0.2cm}		
	\textbf{$\Delta Y(T)$ (coef. $=-0{,}1546$):} even at rarely
	observed levels of the volatility shock, $\pm 0{,}05$
	(roughly the 95\textsuperscript{th} percentile of
	\texttt{$\delta Y(T)$}) up to $\pm 0{,}20$ (close to the
	99.5\textsuperscript{th}--99.9\textsuperscript{th} percentile, near the
	extreme tail of the observed distribution, not a
	typical day) --- the implied predicted shift in tomorrow's return
	is about $-0{,}008\,\%$ to $-0{,}031\,\%$, still smaller than
	typical EUR/USD \emph{bid-ask} spreads and transaction costs on
	any real trading platform. \textbf{This
		reinforces rather than weakens the conclusion}: even a large,
	rarely observed volatility shock --- not just an
	ordinary shock --- does not produce an economically
	meaningful move.
	
	\vspace{0.2cm}
	
	\textbf{$r(t)$ (coef. $=0{,}0206$):} likewise, a typical
	previous-day return of $\pm 0{,}5\,\%$ (roughly the sample
	standard deviation, \S\ref{sec:garch}) would shift tomorrow's
	predicted return by about $\pm 0{,}01\,\%$ --- again, negligible
	economically taken in isolation.
	
	\vspace{0.2cm}
	
	\textbf{In plain terms, for a market observer:} knowing that
	volatility jumped yesterday, or which direction the price moved
	yesterday, provides no exploitable edge --- statistically or
	economically --- for guessing whether EUR/USD will rise or fall
	tomorrow. Anyone building a directional \emph{day-trading}
	rule on either of these two inputs alone
	would, on this evidence, do no better than guessing --- the
	well-known \textbf{\textit{"no \\free lunch"}} conclusion of market efficiency, illustrated here with figures specific to this
	project rather than asserted in the abstract.
\end{econbox}

\begin{figure}[H]
	\centering
	\includegraphics[width=1\linewidth]{04_predicted_effect_with_percentiles.png}
	\caption{Predicted effect on $r(t+1)$ at the $\pm 0{,}05$ / $\pm 0{,}20$ percentiles of $\Delta Y(t)$ --- the magnitude discussed above.}
	\label{fig:baseline-predicted-effect}
\end{figure}

\subsection{What this result does not settle}

\begin{itemize}
	\item Whether $\Delta Y(t)$ explains the \textbf{magnitude} of the return
	($|r(t+1)|$) rather than its direction is not tested at this stage, and
	is a theoretically more natural target given that the GARCH
	itself is a volatility-magnitude model
	(section~\ref{sec:magnitude}).
	\item Whether control variables (rate differential, DXY)
	change this result --- tested next
	(section~\ref{sec:controls}), with a comparison of $\beta_1$
	before/after their addition.
	\item Whether volatility regimes can still be significantly
	detected from $\sigma(t)$ or $\Delta Y(t)$,
	independently of return predictability --- a distinct
	question, addressed at a later stage
	(sections~\ref{sec:regime-hmm} and~\ref{sec:msgarch}).
\end{itemize}

\newpage

% ═════════════════════════════════════════════════════════════════
\section{Regression with controls --- robustness of the baseline}
\label{sec:controls}
% ═════════════════════════════════════════════════════════════════

\subsection{Objective}

Test whether adding two control variables --- the Fed/ECB
interest rate differential and a broad dollar proxy (DXY) --- changes
the baseline result (\S\ref{sec:baseline}: $\beta_1$ on
$\Delta Y(t)$ not significant, model not jointly significant).
The goal is not to seek significance, but to
check whether the "no controls" null result on $\Delta Y(t)$ might have been masking
a confounded relationship.

\subsection{Sourcing the control variables}

The three series come from FRED (the same stable CSV
endpoint already used for EUR/USD)
\begin{table}[H]
	\centering
	\caption{Control variables}
	\begin{tabular}{lll}
		\toprule
		Variable & FRED series & Frequency \\
		\midrule
		Fed Funds rate & \texttt{DFF} & Daily \\
		ECB deposit rate & \texttt{ECBDFR} & Daily \\
		DXY proxy & \texttt{DTWEXBGS} & Daily \\
		\bottomrule
	\end{tabular}
\end{table}

\textbf{Decision: ECB deposit rate, not the MRO or marginal
	lending rate.} FRED also hosts \texttt{ECBMRRFR} (main refinancing
rate) and \texttt{ECBMLFR} (marginal lending facility). The deposit
rate (\texttt{ECBDFR}) is chosen because it is considered the ECB's
\emph{de facto} reference rate under the ample-liquidity
regime under which the Eurosystem has operated for most of
the sample period --- consistent with the recent literature on
FX models based on the rate differential.

\vspace{0.3cm}

\begin{limitbox}
	\texttt{DTWEXBGS} (broad, trade-weighted dollar index)
	is \textbf{not} the commercial DXY index. The DXY weights a basket
	of 6 currencies dominated by the euro ($\approx 58\,\%$ weight);
	\texttt{DTWEXBGS} is a broader basket including emerging
	market currencies. Both measure "dollar strength"
	but are not interchangeable. It was chosen as an accessible
	alternative hosted by FRED, consistent with the project's overall
	sourcing strategy (an imperfect but well-targeted and
	reproducible measure rather than an unreachable proprietary measure ---
	the same logic already applied to the choice of the volatility proxy itself,
	\S\ref{sec:data-sources}).
\end{limitbox}

\subsection{Calendar alignment check}

Before merging, a check for date mismatches between the
control series and the EUR/USD return series: 2 dates present in the
controls but absent from the returns (a start-of-sample edge
effect, and one partial day not recorded by FRED); 20
dates present in the EUR/USD returns but absent from the
controls (mostly identifiable U.S. federal
closures not reflected in \texttt{DFF}, e.g. the
closures linked to Hurricane Sandy in 2012, or winter storm
Jonas in 2016 --- while FX continued to trade via
other financial centers).

\vspace{0.2cm}

\textbf{Conclusion:} 22 misaligned dates out of $\approx 4,175$
($\approx 0{,}5\,\%$), all consistent with identifiable calendar
effects rather than a data quality issue. Handled
via an inner join, not investigated further.

The sample aligned with the controls has 4,153
observations, versus 4,173 in the original baseline
(\S\ref{sec:baseline}) --- a slight reduction due to the calendar
mismatches above. For the before/after comparison of $\beta_1$
to be meaningful, \textbf{the baseline model is
	re-estimated on this same 4,153-observation sample}, rather than
compared directly to the original baseline result on 4,173
observations. This isolates the effect of adding the controls from the
effect of a changed sample.

\subsection{Specification decisions}

Full model:

\begin{keyeq}
	\begin{equation}
		r(t+1) = \alpha + \beta_1 \Delta Y(t) + \beta_2 r(t) + \beta_3\,
		diff_{rate}(t) + \beta_4\, r_{dxy}(t) + u(t)
	\end{equation}
\end{keyeq}

\vspace{0.3cm}
\begin{itemize}
	\item \textbf{	$diff_{rate}(t)$: level, not change.}
	
	
	\begin{center}
		{\large$diff_{rate}(t) = rate_{fed}(t)- rate_{ecb}(t)$}
	\end{center}
	
	
	
	Decision: use the \textbf{level}, not the daily
	change. Policy rates are step functions, only changing on rare
	monetary policy meeting dates --- a change series would be zero on
	almost every day and non-zero only at policy shocks, a fundamentally
	different (event-driven) signal from a continuous uncovered
	interest rate parity-type differential. The level is also the conventional choice in
	the uncovered interest rate parity (UIP) literature, which
	compares rate \emph{levels}, not their changes.
	

\newpage

	\item \textbf{$r_{dxy}(t)$: log return, not level.}
	
	
	\begin{center}
		{\large$r_{dxy}(t) = \ln\left(\frac{dxy_{proxy}(t)}{dxy_{proxy}(t-1)}\right)$}
	\end{center}
	
	
	
	
	
	Decision: use the log return, not the raw index
	level. \texttt{$dxy\_proxy$} is a price-type index, structurally
	similar to EUR/USD itself --- using its level would risk the
	same spurious regression problem (non-stationarity,
	\citep{grangernewbold1974}) already identified and corrected for EUR/USD
	(\S\ref{sec:spec-iterations}). Consistent treatment across all
	price-type series in the model.
\end{itemize}





\subsection{Results}

\begin{table}[H]
	\centering
	\small
	\caption{Baseline re-estimated on the aligned sample ($n=4,153$)}
	\begin{tabular}{lrl}
		\toprule
		& Coef. & Signif. \\
		\midrule
		$\Delta Y(t)$ & $-0{,}1329$ & $p=0{,}578$, not significant \\
		$r(t)$ & $0{,}0196$ & $p=0{,}243$, not significant \\
		\midrule
		\multicolumn{3}{l}{$F=0{,}8512$ ($p=0{,}427$) --- not jointly significant} \\
		\bottomrule
	\end{tabular}
\end{table}

\begin{table}[H]
	\centering
	\small
	\caption{Extended specification with controls ($n=4,153$)}
	\begin{tabular}{lrl}
		\toprule
		& Coef. & Signif. \\
		\midrule
		$\Delta Y(t)$ & $-0{,}1308$ & $p=0{,}584$, not significant \\
		$r(t)$ & $0{,}0235$ & $p=0{,}432$, not significant \\
		$diff_{rate}(t)$ & $0{,}0100$ & $p=0{,}257$, not significant \\
		$r_{dxy}(t)$ & $0{,}8591$ & $p=0{,}865$, not significant (the least of all) \\
		\midrule
		\multicolumn{3}{l}{$F=0{,}7265$ ($p=0{,}574$) --- not jointly significant} \\
		\multicolumn{3}{l}{AIC $=6337$, BIC $=6369$ (vs $6334$ / $6353$ for the re-estimated baseline --- \textbf{both worsen})} \\
		\bottomrule
	\end{tabular}
\end{table}

\vspace{0.4cm}

\begin{table}[H]
	\centering
	\caption{Stability comparison of $\beta_1$ --- the central test of this step}
	\begin{tabular}{lrr}
		\toprule
		& $\beta_1(\Delta Y(t))$ & $p-value$\\
		\midrule
		Before controls & $-0{,}1329$ & $0{,}578$ \\
		After controls & $-0{,}1308$ & $0{,}584$ \\
		\bottomrule
	\end{tabular}
\end{table}

\vspace{0.4cm}

\begin{figure}[H]
	\centering
	\includegraphics[width=1\linewidth]{01_beta1_stability.png}
	\caption{Stability of $\beta_1$ before~/~after adding the controls.}
	\label{fig:beta1-stability}
\end{figure}

$\beta_1$ changes by about $1{,}6\,\%$ in magnitude, the $p-value$ remains
essentially unchanged. \textbf{No evidence of confounding} --- the null
effect of the volatility shock on the direction of the next day's return
is not an artifact of omitting the rate differential or dollar
strength.

\subsection{Interpretation}

This is a \textbf{robustness confirmation of a null result}, not
a new finding. Three points are now established
together:

\begin{enumerate}
	\item The baseline null result (no directional predictability
	from $\Delta Y(t)$ or $r(t)$ alone) does not stem from an
	omitted-variable bias due to the two theoretically most
	obvious FX confounders.
	\item Neither $diff_{rate}$ nor $r_{dxy}(t)$ predict, on their
	own, the direction of the next day's EUR/USD return at this
	horizon --- consistent with the broader empirical literature on
	the difficulty of beating a random walk in FX at short horizons
	\citep{meeserogoff1983}.
	\item Model fit (AIC/BIC) actually \textbf{worsens}
	with the added controls, reinforcing the finding that their
	inclusion is not statistically warranted for this specific
	target variable (next day's return direction) --- a deliberate
	stopping point, not an invitation to keep adding
	variables in search of significance (a
	\emph{p-hacking} risk explicitly avoided, as already noted in
	section~\ref{sec:baseline}).
\end{enumerate}

\newpage
\textbf{What this does not call into question:} as with the baseline,
this null result on \emph{direction} says nothing about whether
$\Delta Y(t)$ explains the \emph{magnitude} of the return
($|r(t+1)|$), nor about the market regime detection
objective --- both remain open and distinct questions
(sections~\ref{sec:magnitude} and~\ref{sec:regime-hmm}).

\subsection{What the two control coefficients mean economically}

\begin{econbox} [$diff_{rate}(t)$ (coef. $=0{,}0100$, $p=0{,}257$)]
	
	Economic theory (uncovered interest rate parity, UIP)
	predicts a relationship between the rate differential and the
	expected exchange return. Taking this coefficient at face value despite its lack
	of significance: a one-percentage-point widening of the
	Fed-ECB rate differential (a fairly large policy move,
	larger than most individual rate decisions) would be
	associated with a shift of only $0{,}01\,\%$ in the expected
	next-day return --- an effect too small to matter
	economically, even before considering that it is statistically
	indistinguishable from zero. This null result is not an anomaly
	specific to this project: it echoes a well-known puzzle in
	international finance (the "\emph{forward premium puzzle}",
	\citep{fama1984}), where UIP frequently fails to hold, or
	holds with an unexpected sign, in short-horizon empirical
	tests. This project's finding is consistent with this broader
	literature, decades old, rather than contradicting it.
\end{econbox}

\vspace{2cm}

\begin{econbox}[$r_{dxy}(t)$ (coef. $=0{,}8591$, $p=0{,}865$)]
	
	The coefficient's 95\,\% confidence interval runs from about
	$-9{,}0$ to $+10{,}7$. In plain terms, this range is compatible with
	either "the dollar index and the next day's EUR/USD returns
	move strongly in opposite directions" or
	"they move strongly in the same direction" or
	"there is no relationship at all" --- the data at this
	daily horizon simply do not pin down the relationship one way
	or the other. This is a materially different (and more
	informative) finding than "not significant" alone: it says that
	the estimate carries essentially no exploitable precision,
	not merely weak or moderate precision.
\end{econbox}

\begin{center}
	\begin{figure}[H]
		\centering
		\includegraphics[width=1\linewidth]{04_r_dxy_wide_ci.png}
		\caption{95\,\% confidence interval of the $r_{dxy}(t)$ coefficient, $[-9{,}0~;~+10{,}7]$.}
		\label{fig:r-dxy-wide-ci}
	\end{figure}
\end{center}


\vspace{0.8cm}

\begin{econbox}[Practical conclusion for anyone building on this model]
	At a one-day horizon, neither the rate differential nor
	broad dollar strength provides an exploitable directional signal for
	EUR/USD, beyond what is already known (nothing) from the volatility
	shock and momentum terms alone. This does not mean
	these variables are economically irrelevant to exchange rates
	in general --- the literature is clear that they matter at
	other horizons and in other specifications --- only that,
	in this precise specification (level/return, one day
	ahead), they add no exploitable information.
\end{econbox}

\newpage

\subsection{Robustness check of the standard errors, motivated by the strong autocorrelation of "$diff_{rate}(t)$"}

$diff_{rate}(t)$ (the level of the Fed-ECB rate differential) is
extremely persistent --- autocorrelation of $0{,}998$ at lag~1,
still $0{,}98$ at lag~20 (an expected property of a policy-rate
level series, which moves in rare discrete steps and stays
flat in between; see the same point about its step-function
nature, \S\ref{sec:controls}). This raised a concern: the
HAC (Newey-West) standard errors reported above used
\texttt{maxlags=5}, potentially insufficient to fully correct
for such persistent autocorrelation --- an under-corrected
standard error would be too small, potentially overstating
significance.

\textbf{Check performed:} re-running the extended
specification with \texttt{maxlags=10} and \texttt{maxlags=20}. The result does not
change in any meaningful way: \texttt{diff\_rate} and all the
other coefficients remain far from conventional
significance thresholds, regardless of the lag choice. This is the
expected result, not a surprising one --- an overly short
\texttt{maxlags=5} would, if anything, have \emph{underestimated} the
true standard errors, meaning the already non-significant result would only become
\emph{more} clearly non-significant with a longer lag
window, never less. The robustness check confirms this
rather than calling the conclusions above into question.

\newpage

% ═════════════════════════════════════════════════════════════════
\section{Magnitude regression --- does the volatility \\ signal  explain $|r(t+1)|$?}
\label{sec:magnitude}
% ═════════════════════════════════════════════════════════════════

\subsection{Objective and how this test differs from the previous two}

The baseline and control regressions without controls (section~\ref{sec:baseline}) and with controls
(section~\ref{sec:controls}) tested whether the volatility shock explains the
\textbf{direction} of the next day's return --- $r(t+1)$, a
signed quantity. Both returned a robust null result.

This step tests a different target, theoretically more natural
for a volatility model: the \textbf{magnitude} of the next day's
return --- $|r(t+1)|$, the size of the move regardless of its
sign.

\textbf{Important distinction to keep in mind throughout
	this section:} $|r(t+1)|$ carries \textbf{no directional
	information}. Whether the actual return is $+0{,}8\,\%$ or
$-0{,}8\,\%$, $|r(t+1)| = 0{,}8\,\%$ in both cases. A model that
explains $|r(t+1)|$ well \textbf{is not} a trading signal saying
whether to buy or sell --- it is a signal about the size of the move to
expect, in either direction. This is exactly the
"price range" framing of the project's founding question
(\S\ref{sec:intro}): not a point forecast, but an interval that
widens or narrows with the volatility regime.

\subsection{Why a significant result here would not be a new discovery}

This point was raised and clarified before any code was run, and
is kept here in full because it conditions how the results
below should be read.

The GARCH model is defined by construction as:

\begin{keyeq}
	\begin{align}
		r(t+1) &= \mu + \sigma(t+1) \cdot z(t+1), \qquad z(t+1) \sim \text{i.i.d.}(0,1) \\
		\sigma^2(t) &= \omega + \alpha\,\varepsilon^2(t-1) + \beta\,\sigma^2(t-1)
	\end{align}
\end{keyeq}

\vspace{0.3cm}

In absolute value: $|r(t+1)| \approx \sigma(t+1)\cdot|z(t+1)|$. And
$\sigma(t+1)$ is \textbf{itself built from $\sigma(t)$}
via the GARCH recursion --- with $\beta \approx 0{,}961$ (highly
persistent, confirmed in section~\ref{sec:garch}). 

Asking whether "$\sigma(t)$
explains $|r(t+1)|$?" therefore comes close to asking
"did the GARCH do what it was calibrated to do?" --- a
question already answered affirmatively by the diagnostics of
section~\ref{sec:garch} (Ljung-Box on the squared standardized
residuals, not significant $\Rightarrow$ the model correctly captures
magnitude dynamics).

\vspace{0.1cm}

\textbf{Consequence for interpretation, decided in advance:} a
significant result here should be documented as an
\textbf{internal consistency check of the GARCH}, not as a
new empirical discovery about FX behavior --- unlike
the direction test, where the null result was itself the finding,
consistent with, and adding a data point to, the weak-form
market efficiency literature.

\vspace{0.1cm}

\textbf{The mirror-image expectation, also decided in advance:} a
\emph{null} result here would be the surprising result, and would point
to a likely specification problem rather than a genuine
discovery that "magnitude too is unpredictable" --- the reverse
logic of the direction test.

\subsection{Two specifications tested side by side}

\textbf{The dilution concern.} 

\vspace{0.3cm}

Throughout this project,

\begin{center}
	
	
	{\large $\Delta Y(t) = \ln\left(\frac{\sigma(t)}{\sigma(t-1)}\right)$}
\end{center}

\vspace{0.3cm}

The log change in conditional volatility --- has been the standard "shock" variable
used for the level-shift signal. But
$\sigma(t)$ is strongly persistent (autocorrelated by construction).
Taking its log difference could dilute a strong, obvious relationship
at the \textbf{level}, since a daily change captures much less
information than the level itself when the underlying series moves
little, day after day, relative to its own scale.

To distinguish "the magnitude relationship is genuinely weak" from
"the log-change transformation dilutes an otherwise strong
level relationship," both are tested:

\vspace{0.3cm}


\begin{keyeq}
	
	\begin{itemize}
		\item \textbf{Spec A:} $|r(t+1)| = a + b_1 \cdot \Delta Y(t) + b_2 \cdot |r(t)| + u(t)$
		\item \textbf{Spec B:} $|r(t+1)| = a + b_1 \cdot \sigma(t)~\text{(level)} + b_2 \cdot |r(t)| + u(t)$
	\end{itemize}
\end{keyeq}
\vspace{0.3cm}

\textbf{The control variable also changes form: $|r(t)|$, not
	$r(t)$.} In the direction regressions, $r(t)$ tested
momentum~/~mean reversion. That is not the relevant question
for magnitude. What matters here is whether the \emph{size} of
yesterday's move predicts the size of tomorrow's move --- another
facet of volatility \emph{clustering}. The control is therefore
switched to $|r(t)|$, decided to be \textbf{mandatory} (not optional)
at this step, given the change in what is being explained.

\newpage

\subsection{Results}

Sample size: 4,173 observations (same alignment as the
original baseline).

\begin{table}[H]
	\centering
	\small
	\caption{Spec A --- $\Delta Y(t)$}
	\begin{tabular}{lrrrr}
		\toprule
		& Coef. & Std. err. & $z$ & $P>|z|$ \\
		\midrule
		const & $0{,}3386$ & $0{,}009$ & $39{,}499$ & $0{,}000$ \\
		$\Delta Y(t)$ & $-0{,}0871$ & $0{,}159$ & $-0{,}549$ & $0{,}583$ --- \textbf{not significant} \\
		$|r(t+1)|$ & $0{,}1148$ & $0{,}017$ & $6{,}562$ & $0{,}000$ --- significant \\
		\midrule
		\multicolumn{5}{l}{$R^2 = 0{,}013$ ; $F=22{,}04$ ($p=3{,}01\times10^{-10}$, model jointly significant)} \\
		\multicolumn{5}{l}{AIC $=3026{,}94$ ; BIC $=3045{,}95$} \\
		\bottomrule
	\end{tabular}
\end{table}

The model is jointly significant, but \textbf{only via
	$|r(t)|$}, not via $\Delta Y(t)$. $\Delta Y(t)$ itself
carries no detectable explanatory power here.

\vspace{0.1cm}

\begin{table}[H]
	\centering
	\small
	\caption{Spec B --- $\sigma(t)$ level}
	\begin{tabular}{lrrrr}
		\toprule
		& Coef. & Std. err. & $z$ & $P>|z|$ \\
		\midrule
		const & $0{,}0171$ & $0{,}018$ & $0{,}946$ & $0{,}344$ \\
		$\sigma(t)$ & $0{,}7029$ & $0{,}040$ & $17{,}364$ & $0{,}000$ --- \textbf{strongly significant} \\
		$|r(t+1)|$ & $0{,}0269$ & $0{,}016$ & $1{,}638$ & $0{,}101$ --- not significant here \\
		\midrule
		\multicolumn{5}{l}{$R^2 = 0{,}090$ ; $F=174{,}1$ ($p=2{,}52\times10^{-73}$)} \\
		\multicolumn{5}{l}{AIC $=2687{,}29$ ; BIC $=2706{,}30$} \\
		\bottomrule
	\end{tabular}
\end{table}

\begin{figure}[H]
	\centering
	\includegraphics[width=1\linewidth]{02_specB_scatter_sigma.png}
	\caption{Spec B --- scatter plot of $|r(t+1)|$ against $\sigma(t)$, the only positive relationship in the whole project.}
	\label{fig:specB-scatter}
\end{figure}

\vspace{0.1cm}

$\sigma(t)$ is strongly significant. Once included, $|r(t)|$
loses its significance --- consistent with $\sigma(t)$
already absorbing the persistence information that $|r(t)|$ alone
approximated in Spec~A.

\vspace{0.1cm}

\begin{table}[H]
	\centering
	\caption{Spec A / Spec B comparison}
	\begin{tabular}{lrr}
		\toprule
		& Spec A ($\Delta Y(t)$) & Spec B ($\sigma(t)$ level) \\
		\midrule
		$R^2$ & 0.013 & \textbf{0.090} \\
		AIC & 3026.94 & \textbf{2687.29} \\
		BIC & 3045.95 & \textbf{2706.30} \\
		Shock variable significant? & No ($p=0{,}583$) & \textbf{Yes ($p<0{,}001$)} \\
		\bottomrule
	\end{tabular}
\end{table}

\vspace{0.4cm}

\begin{figure}[H]
	\centering
	\begin{subfigure}{0.48\linewidth}
		\includegraphics[width=\linewidth]{03_specA_vs_specB_fit.png}
		\caption{Fit comparison, Spec A vs Spec B}
	\end{subfigure}
	\hfill
	\begin{subfigure}{0.48\linewidth}
		\includegraphics[width=\linewidth]{04_significance_flip.png}
		\caption{Significance flip of the shock variable}
	\end{subfigure}
	\caption{The two closing figures of the results section: fit comparison, then significance flip.}
	\label{fig:specA-vs-specB}
\end{figure}

\vspace{0.2cm}

\textbf{Conclusion on the dilution question:} confirmed. Spec~B
dominates Spec~A on every criterion ($R^2$, AIC, BIC). The
log-change transformation used elsewhere in the project
significantly weakens the signal here --- \textbf{\textit{it is the level of
		conditional volatility, not its daily change, that carries
		explanatory power for the magnitude of the next day's return}}.

\subsection{What an $R^2$ of 0.09 means economically --- a worked example}

\begin{econbox}[$R^2$ of 0.09]
	Using the fitted Spec~B coefficients (const $=0{,}0171$,
	sigma\_t coefficient $=0{,}7029$):
	\begin{itemize}
		\item $\sigma(t) \approx 0{,}25$ (corresponding roughly
		to the \textbf{0.1\textsuperscript{th} percentile} of the observed
		$\sigma(t)$ distribution --- essentially the practical floor of the
		sample, not merely "relatively low"):
		predicted $|r(t+1)| \approx 0{,}0171 + 0{,}7029\times0{,}25
		\approx \mathbf{0{,}19\,\%}$.
		\item $\sigma(t) \approx 0{,}60$ (corresponding roughly
		to the \textbf{71.6\textsuperscript{th} percentile} --- just above
		the median, not a turbulent or post-shock extreme as
		the initial framing suggested):
		predicted $|r(t+1)| \approx 0{,}0171 + 0{,}7029\times0{,}60
		\approx \mathbf{0{,}44\,\%}$.
	\end{itemize}
\end{econbox}

\newpage	

\textbf{\textit{NB}}:


These two values are not two symmetric, representative points
of "calm vs turbulent" --- one sits at the extreme floor of the
distribution, the other just above its
center. Calling 0.60 "turbulent" or
"shortly after a shock" overstated how unusual this level
actually is. The qualitative finding survives (a higher
$\sigma(t)$ does imply a wider expected move), but the
"roughly twice as wide" comparison above should be read
as floor-versus-median, not calm-versus-crisis. A properly
matched pair --- for example the 10\textsuperscript{th}
and 90\textsuperscript{th} percentiles of $\sigma(t)$ --- would better
illustrate a genuine calm-versus-turbulent contrast; producing
this pair requires re-reading the percentiles directly from the
processed data, not done here.


In plain terms: even comparing the practical floor of volatility to a
fairly ordinary, near-median level, the model's expected move
roughly doubles --- with zero information about which direction
that move will take. This is directly usable for
the project's operational objective --- a volatility-adjusted
price range~/~risk sizing signal
(section~\ref{sec:price-range}) --- even though it says nothing about
direction.

\vspace{0.2cm}

$R^2=0{,}09$ --- is this "a lot"? Only relative to the
general difficulty of FX forecasting --- the direction
regressions returned $R^2 \approx 0{,}000$--$0{,}001$, so 0.09
represents a genuine, usable step forward in explanatory power
in this context. In absolute terms, $91\,\%$ of the variance of
move size remains unexplained; this is a modest but
real signal, not a breakthrough.

\subsection{Was predictive power detected? Yes --- but for a narrower claim than "prediction"}

\begin{table}[H]
	\centering
	\small
	\caption{Direction vs magnitude --- summary}
	\begin{tabular}{|p{4.7cm}|p{4.2cm}|p{4.2cm}|}
		\hline
		& Direction --- $r(t+1)$ & Magnitude --- $|r(t+1)|$ \\
		\hline
		Predictive power detected? & No (robust null) & \textbf{Yes} ($R^2=0{,}090$, $p<0{,}001$) \\
		\hline
		Usable as a directional bet? & No --- confirmed unusable & No --- magnitude never carries direction \\
		\hline
		Usable as an adjusted range? & --- & \textbf{Yes} --- exactly this signal \\
		\hline
		Scientifically surprising finding? & Somewhat --- a well-documented robust null is informative in itself & No --- expected almost directly from the GARCH's own construction (\S~above) \\
		\hline
	\end{tabular}
\end{table}

Both findings are true at the same time, without contradiction: the
magnitude result is a \textbf{real,
	statistically robust, and operationally exploitable} finding, and at
the same time \textbf{not a new scientific discovery}, since it
follows closely from how the GARCH itself is built. Neither
framing cancels out the other.

\newpage

% ═════════════════════════════════════════════════════════════════
\section{Regime detection via a Gaussian HMM}
\label{sec:regime-hmm}
% ═════════════════════════════════════════════════════════════════

\subsection{Objective and conceptual framing decided before coding}

Objective (b) of the project, stated from the outset, was to detect
market regimes described as "low~/~normal~/~high."
Before any implementation, this framing was revisited and revised.

%	\textbf{Why "bubble" is not an accurate label for what this model detects.} The build-up phase of a financial bubble is often characterized by \textbf{low and stable} volatility --- a steady, discreet price appreciation is part of what makes a bubble appear "safe," and attracts more capital. It is typically the \textbf{burst}, not the bubble itself, that produces a volatility spike. A purely volatility-based model (this HMM, or the underlying GARCH) has no way to distinguish an "ordinary calm market" from a "silently forming bubble" --- both show an identical low-volatility signature. Labeling a detected low-volatility state as "no bubble" and a high-volatility state as "bubble" would be an unwarranted conflation of two distinct economic concepts: a volatility regime (directly measurable here) and a valuation regime (would require a measure of price deviation from fundamentals, outside the scope of this project).

%	\textbf{Decision:} states are labeled according to what the model
%	actually measures --- \textbf{"low volatility" /
	%		"normal volatility" / "high volatility
	%		(stress~/~panic)"} --- not "normal~/~bubble~/~panic".
%	This revises the original framing from the introduction for accuracy,
%	rather than assuming the initial wording was correct once
%	implementation raised the question.

\subsection{Input variable and transformation}

\textbf{Input:} $\sigma(t)$, the conditional volatility of the
GARCH(1,1) (Student's~$t$ specification, section~\ref{sec:garch}) ---
the level, not its change $\Delta Y(t)$, consistent with the
finding in section~\ref{sec:magnitude} that the level carries
noticeably more explanatory information than the log change.

\textbf{Transformation applied:} $\log(\sigma(t))$, not raw
$\sigma(t)$. This is a \textbf{level} transformation (monotone, preserves
order) --- not to be confused with \\ $\Delta Y(t) =
\ln(\sigma(t)/\sigma(t-1))$, which is a \textbf{change} (a
genuinely different quantity). The log transform is applied
because \texttt{GaussianHMM} assumes each hidden state generates
approximately normally distributed observations; $\sigma(t)$
is strictly positive and likely right-skewed in its
raw form, a poor fit to this assumption --- the same
reasoning already applied when moving from EUR/USD's raw price
level to log returns (\S\ref{sec:spec-iterations}).



\subsection{Results}

\textbf{Convergence check:} 8 of 10 restarts converge
to the identical log-likelihood (3302.55) --- strong evidence that this
is the \textbf{global optimum}, not a new local optimum. %This resolves the previous section's issue --- it was mainly an initialization problem, not a fundamental inadequacy of a Gaussian HMM for this series (the structural concern above still bears keeping in mind for future extensions, but did not materialize as a blocking problem here).

\begin{table}[H]
	\centering
	\caption{Regime summary}
	\begin{tabular}{llrrrr}
		\toprule
		State & Label & Mean($\log\sigma$) & Std. dev. & Frequency & Approx. $\sigma$ \\
		\midrule
		0 & Low volatility & $-1{,}0725$ & 0.1268 & $29{,}0\,\%$ & 0.342 \\
		1 & Normal volatility & $-0{,}7476$ & 0.0885 & $38{,}1\,\%$ & 0.474 \\
		2 & High volatility (stress~/~panic) & $-0{,}3907$ & 0.1081 & $32{,}9\,\%$ & 0.677 \\
		\bottomrule
	\end{tabular}
\end{table}

The three states are now clearly separated (means
$-1{,}07$~/~$-0{,}75$~/~$-0{,}39$, non-overlapping given
their standard deviations), the opposite of the first degenerate
attempt.

\begin{figure}[H]
	\centering
	\includegraphics[width=1\linewidth]{01_sigma_by_regime_fullsample.png}
	\caption{$\sigma(t)$ by regime, full sample --- the classification the table above describes in figures.}
	\label{fig:sigma-by-regime}
\end{figure}

\textbf{Transition matrix and implied regime persistence.}
The diagonal (staying in the same regime) is very high for the
low- and high-volatility states (0.9921 and 0.9910), consistent
with genuine persistence rather than daily noise. Using the
standard expected regime-duration formula, $1/(1-p_{\text{stay}})$:

\vspace{0.3cm}

\begin{itemize}
	\item \textbf{Low volatility:} $1/(1-0{,}9921) \approx
	\mathbf{127}$ trading days ($\approx$~6 months)
	\item \textbf{High volatility (stress~/~panic):}
	$1/(1-0{,}9910) \approx \mathbf{111}$ days ($\approx$~5 months)
	\item \textbf{Normal volatility} ($p_{\text{stay}} \approx 0{,}9861$):
	$1/(1-0{,}9861) \approx \mathbf{72}$ days ($\approx$~3.5 months)
\end{itemize}

\vspace{0.3cm}

These durations are of the same order of magnitude as some of the
half-lives found in section~\ref{sec:garch-stability} (e.g.\ Periods 1, 4, 5 in
the 65--80-day range) --- an approximate, but sensible,
consistency check between the two independent analyses, without
constituting a formal validation.

\vspace{0.2cm}

\textbf{Most recent observations.} The last 10 trading
days of the sample (2026-08-28 to 2026-09-11) are all classified as
"low volatility" --- a stable, single-regime segment, not
a daily flip-flop as in the first (degenerate) attempt.

\newpage

\subsection{Why Gaussian emissions here, despite the Student's $t$ preference for the GARCH itself}

\textbf{They model different objects.} The GARCH's
Student's~$t$ choice concerns the distribution of return
innovations $z(t)$, which show genuine fat tails (empirical excess
kurtosis $\approx 1{,}83$, section~\ref{sec:garch}). This HMM instead
models $\log(\sigma(t))$ \textbf{conditional on a regime}.
Once observations are split by regime, the residual
distribution \emph{within} a single regime can be reasonably
close to Gaussian --- because genuine extreme events are
absorbed by the \textbf{regime-switching structure itself}
(they are assigned to the "high volatility" state) rather than
remaining outliers within a single state's distribution. A
mixture of several Gaussians (one per regime) can
collectively reproduce a fat-tailed \emph{overall}
distribution even if each individual component is Gaussian --- the
fat tail is captured by the regime structure, not by
within-regime kurtosis.

\textbf{A practical constraint, also stated plainly:}
\texttt{hmmlearn}, the library used here, does not natively
support Student's~$t$ emissions --- only Gaussian or a
Gaussian mixture (\texttt{GMMHMM}). Implementing
Student's~$t$ emissions would require a different, less mature
library (e.g.\ \texttt{pomegranate}) or a custom-coded HMM, a
level of complexity judged disproportionate here given that the
results above show the Gaussian specification converges cleanly and produces
well-separated, persistent regimes.

\subsection{What these results mean economically}

\begin{econbox}
	
	The market does not spend most of its time in a single "normal" state.
	Over 2010--2026, the sample splits roughly into
	$29\,\%$ low volatility, $38\,\%$ normal volatility, $33\,\%$
	high volatility. In other words, the EUR/USD market spends about
	\textbf{one third of all trading days} in a
	stress~/~panic-level volatility regime --- this is not a rare
	tail event confined to a few weeks of crisis; it is a
	recurring state, roughly one day in three, over a 16-year span.
	Economically, this cuts against a framing in which "panic"
	would be treated as an exceptional, fleeting deviation from a
	normal baseline --- here, it is closer to a third mode of
	operation that the market regularly passes through (wars, central bank
	shocks, COVID, inflation shocks, etc., recurring roughly
	every few years in this sample).
\end{econbox}

\begin{figure}[H]
	\centering
	\includegraphics[width=1\linewidth]{04_regime_frequency.png}
	\caption{Frequency of the three regimes over 2010--2026 --- 29\,\% / 38\,\% / 33\,\%.}
	\label{fig:regime-frequency}
\end{figure}

\vspace{0.9cm}

\begin{econbox}
	Regimes evolve slowly; this is not daily noise.
	
	The implied durations (previous section) --- $\approx 127$~days
	low vol., $\approx 111$~days high vol., $\approx 72$~days
	normal --- mean, in economic terms, that \textbf{once
		EUR/USD enters a volatility regime, it tends to stay there
		for several months}, not days or weeks. This matters
	operationally: a regime signal from this model is not
	a short-term trading trigger --- it is closer to a
	medium-term risk-posture indicator (useful, for instance,
	for sizing positions or a hedging policy over a
	quarter, not for daily decisions).
\end{econbox}

\begin{figure}[H]
	\centering
	\includegraphics[width=1\linewidth]{06_regime_implied_duration.png}
	\caption{Average implied duration of each regime --- 72 to 127 trading days depending on the regime.}
	\label{fig:regime-duration}
\end{figure}

\vspace{0.9cm}
\begin{center}
	\begin{econbox} [What the "current" classification means in practice]
		
		At the end of the sample (2026-09-11), the model classifies the market
		as being in a \textbf{low volatility} regime, and --- according to the
		transition matrix --- regimes of this type persist, on
		average, more than 4 months once entered. In plain terms: based on
		this signal alone, the market has recently been calm and,
		historically, calm regimes tend to persist rather than reverse
		abruptly. This is explicitly \textbf{not} a
		forecast that the market will remain calm (the model has no
		crystal ball on the \emph{timing} of the next regime
		change), and it says nothing about the future
		direction of EUR/USD (cf.\ the direction~/~magnitude
		distinction, section~\ref{sec:magnitude}) --- only
		that, historically, the current regime type does not tend to flip
		overnight.
	\end{econbox}
	
\end{center}

\newpage

\subsection{Known limitations}

\begin{limitbox}
	\begin{itemize}
		\item \textbf{Inherited GARCH instability} (documented in
		section~\ref{sec:garch-stability}): the $\sigma(t)$ used as
		input here comes from a single GARCH fit on the full
		2010--2026 sample, already shown to have persistence
		that varies between sub-periods (though not statistically
		established, \S\ref{sec:garch-stability}). The regimes
		detected by this HMM may partly reflect this
		averaging artifact rather than purely genuine market regime
		shifts. This is the documented justification for treating this HMM as a
		first-approximation approach, with a Markov-Switching
		GARCH (re-estimating GARCH parameters by regime, directly
		addressing the instability) flagged as a possible refinement
		(section~\ref{sec:msgarch}).
		\item \textbf{"High volatility" is not "bubble"}
		--- this limitation is structural to any purely
		volatility-based regime model, not specific to this
		implementation, and should not be silently reinterpreted
		as such in a later summary of the project.
		\item Regime labels and boundaries depend on the arbitrary but
		standard choice $n\_states=3$; not tested against 2 or 4
		states as alternatives at this stage.
	\end{itemize}
\end{limitbox}

\newpage

% ═════════════════════════════════════════════════════════════════
\section{Markov-Switching GARCH ($K=2$)}
\label{sec:msgarch}
% ═════════════════════════════════════════════════════════════════

This model is motivated by the sub-period variation in
persistence reported in section~\ref{sec:garch-stability} (variation in the
point estimates; that section's formal likelihood-ratio test does
not establish instability, $p=0{,}277$) and addresses it
more directly than the "HMM on a fixed GARCH" approach of
section~\ref{sec:regime-hmm}.

\subsection{Why R, not Python --- a documented and \\deliberate exception}

Every previous step of this project was implemented in Python
(\texttt{arch} for the GARCH, \texttt{hmmlearn} for the HMM). The
Markov-Switching GARCH is different: no comparably
mature Python package exists. The standard reference is the R package
"\textbf{MSGARCH}"
\citep{ardia2019}, built in C++/Rcpp.

This is not merely a lack of tooling. Estimating a true
MS-GARCH faces a genuine theoretical difficulty, documented since
\citet{hamiltonsusmel1994}: in the fully general model, the
conditional variance of a given regime depends, in principle, on
the entire history of regime paths since the start of the series
(because $\sigma^2(t)$ depends on $\sigma^2(t-1)$, itself depending on
the active regime at $t-1$, and so on) --- making the exact
likelihood of this fully general model intractable beyond a
handful of observations, since it requires summing over an
exponentially growing number of regime paths.

Two different responses exist in the literature.
\citet{gray1996} keeps the fully general model and makes it
\emph{tractable by approximating it}: at each step, the $k$
regime-conditional variances are collapsed into a single
expected variance (weighted by filtered regime probabilities),
used as a common lagged-variance input going forward --- a genuine
approximation of the fully general model's likelihood.

The MSGARCH we implemented does not do this. It implements the \\specification of
\citet{haasmittnikpaolella2004}, which \emph{sidesteps} the
path-dependency problem rather than approximating it: each regime
$k$ runs its own self-contained GARCH(1,1) recursion
($\sigma_k^2(t) = \omega_k + \alpha_k\, r(t-1)^2 + \beta_k\,
\sigma_k^2(t-1)$), fed by the same observed return but its
own lagged variance --- no cross-regime collapsing at any
step. The density of the observed return is then the $K$-regime
mixture of these $K$ parallel processes, weighted by
regime probabilities. This gives an \emph{exact} likelihood
for this (more restricted) model --- it is not an approximation in
Gray's sense. The real cost is different: each regime's
path ignores the history of the regime that was \emph{actually}
realized at each past date --- a simplification relative to
the fully general model, but not the same type of approximation as
Gray's collapsing.


\begin{figure}[H]
	\centering
	\includegraphics[width=1\linewidth]{02_conditional_volatility_by_regime.png}
	\caption{$\sigma(t)$ by regime, full sample --- two independent recursions, never mixed.}
	\label{fig:cond-vol-by-regime}
\end{figure}

Given this, and the availability of R, this step was carried out
as a \textbf{single, deliberate, documented exception} to an
otherwise all-Python project: an R script (\texttt{ms\_garch.R})
run directly (not via \texttt{rpy2}, to avoid the overhead of a
Python-R bridge for a one-off analysis), producing CSV outputs
reloaded into the Python project for further work.

\subsection{Choice of $K=2$, not $K=3$}

The HMM in section~\ref{sec:regime-hmm} used 3 states
(low~/~normal~/~high volatility). This step deliberately
uses \textbf{$K=2$}, not 3, for a stated, not
arbitrary, reason: estimating an MS-GARCH is substantially more
numerically difficult than fitting a Gaussian HMM on an
already-computed scalar series ($\sigma(t)$) --- it must jointly
estimate $\omega, \alpha, \beta$ \textbf{by regime}, plus the
transition matrix, under the parallel-regime specification of
Haas et al.\ (2004a) \citet{haasmittnikpaolella2004} described above. Each additional regime
multiplies this difficulty. 

The FX~/~equity MS-GARCH literature
(e.g.\ \citet{klaassen2002}; \citet{marcucci2005}) usually
starts with 2 regimes before considering more. This is a
genuine methodological difference from the HMM, not an
inconsistency --- the two tools face different numerical
constraints.

\newpage

\subsection{Specification and errors encountered}

\textbf{Model specification.} 2 regimes, each a GARCH(1,1)
("sGARCH") with Student's~$t$ innovations ("std") ---
matching the single-regime specification already validated
(section~\ref{sec:garch}), with one difference: no mean
equation is estimated here (returns are treated as zero-mean),
whereas the single-regime \texttt{arch} fit estimates a
constant mean $\mu$.



\subsection{Results}

\begin{table}[H]
	\centering
	\small
	\caption{Fitted parameters, MS-GARCH ($K=2$)}
	\begin{tabular}{lrrrl}
		\toprule
		& Estimate & Std. err. & $t$ & $\Pr(>|t|)$ \\
		\midrule
		$\omega_1$ (Regime 1) & 0.0001241 & 0.0002112 & 0.587 & 0.278 \\
		$\alpha_1$ (Regime 1) & 0.0106220 & 0.0077169 & 1.376 & 0.084 \\
		$\beta_1$ (Regime 1) & 0.9864591 & 0.0013367 & 737.970 & $<10^{-16}$ \\
		$\nu_1$ & 7.0711 & 0.9083 & 7.785 & $<10^{-14}$ \\
		$\omega_2$ (Regime 2) & 0.0057432 & 0.0039894 & 1.440 & 0.075 \\
		$\alpha_2$ (Regime 2) & 0.0090000 & 0.0132647 & 0.678 & 0.249 \\
		$\beta_2$ (Regime 2) & 0.9802132 & 0.0083909 & 116.818 & $<10^{-16}$ \\
		$\nu_2$ & 22.0046 & 16.8510 & 1.306 & 0.096 (imprecise) \\
		$P_{1\to1}$ & 0.99753 & 0.00311 & 320.741 & $<10^{-16}$ \\
		$P_{2\to1}$ & 0.00627 & 0.00130 & 4.809 & $<10^{-6}$ \\
		\midrule
		\multicolumn{5}{l}{LL $=-2806{,}64$ ; AIC $=5633{,}27$ ; BIC $=5696{,}64$} \\
		\multicolumn{5}{l}{Stable probabilities: Regime~1 $=71{,}7\,\%$, Regime~2 $=28{,}3\,\%$} \\
		\bottomrule
	\end{tabular}
\end{table}

The model has 10 free parameters (four per regime, plus the two
free transition probabilities). The full transition matrix:

\vspace{0.3cm}

\begin{table}[H]
	\centering
	\caption{Transition matrix}
	\begin{tabular}{lrr}
		\toprule
		From \textbackslash{} To & Regime~1 & Regime~2 \\
		\midrule
		Regime~1 & 0.99753 & 0.00247 \\
		Regime~2 & 0.00627 & \textbf{0.99373} \\
		\bottomrule
	\end{tabular}
\end{table}

\vspace{0.3cm}

\textbf{Which regime is "high volatility"?} Long-run
(unconditional) variance by regime, $\omega/(1-\alpha-\beta)$

\begin{table}[H]
	\centering
	\caption{Persistence and long-run $\sigma$ by regime}
	\begin{tabular}{lrrr}
		\toprule
		Regime & $\alpha+\beta$ & Half-life (within-regime persistence) & Implied long-run $\sigma$ \\
		\midrule
		1 & 0.9971 & 237.1 days & 0.206 \\
		2 & 0.9892 & 63.9 days & 0.730 \\
		\bottomrule
	\end{tabular}
\end{table}


\textbf{Regime~2 is the high-volatility regime}, occurring
$28{,}3\,\%$ of the time: on days when its smoothed probability exceeds
$50\,\%$ (1,080 days), the realized daily standard deviation of
returns is $0{,}70$, versus $0{,}44$ on the remaining 3,095
days --- a ratio of $1{,}6$ (the classification is itself based
on volatility, so this is an indicative ratio, not an
independent test). The implied long-run $\sigma$ from the parameters
($0{,}73$ vs $0{,}21$, ratio of $3{,}5$) is a much
wider, much less precise figure: $\omega$ is not significant in either
regime ($p=0{,}278$ and $0{,}075$) and $\alpha+\beta$ is close to
1, so the denominator $1-\alpha-\beta$ is small ($0{,}0029$
and $0{,}0108$) and poorly determined. It should not be read as a
measure of realized volatility.

\vspace{0.3cm}

\begin{limitbox}
	The HMM in section~\ref{sec:regime-hmm} has three states, so
	its "high volatility" state ($32{,}9\,\%$ of the time) and this
	model's Regime~2 do not partition the sample the same
	way: the closeness of $32{,}9\,\%$ and $28{,}3\,\%$ should be read
	as order-of-magnitude consistency, not as an
	independent confirmation.
\end{limitbox}

\vspace{0.3cm} 

\subsubsection{A counter-intuitive finding --- suggestive, not established}

In the point estimates, the high-volatility regime has
\textbf{lower} GARCH persistence ($\alpha+\beta=0{,}989$, half-life
63.9 days) than the low-volatility regime ($\alpha+\beta=0{,}997$,
half-life 237.1 days). In plain terms, the point estimates say that
\textbf{acute stress fades faster than calm periods last}.

\vspace{0.6cm}

\begin{limitbox}
	This should be read as suggestive, not as an established result.
	The $\alpha+\beta$ gap is 0.0079. Using the standard errors
	above and ignoring the covariance between $\alpha$ and $\beta$
	within each regime (not saved with the estimates),
	the standard error of this gap is about 0.018, giving $t \approx
	0{,}45$. The covariances would refine this figure, but nothing here
	supports a statistically significant difference. Moreover, a
	half-life $\ln(0{,}5)/\ln(\alpha+\beta)$ is extremely sensitive
	when $\alpha+\beta$ is this close to 1 (see the Period~2
	diagnosis, \S\ref{sec:garch-stability-p2}).
\end{limitbox}

If the gap is real, it makes economic sense --- a crisis
shock (e.g.\ a central bank surprise, a liquidity panic) is
intense but tends to fade within a few months, whereas a
calm regime, once established, can persist for almost a full year before
the next disturbance.

\newpage

\textbf{Important distinction, easy to conflate:} this
within-regime half-life (how fast GARCH shocks dissipate \emph{within}
a regime) is a different concept from the Markov chain's own
regime duration (how long the market \emph{stays} in a
given regime before switching). Using
$1/(1-P)$ with the transition matrix above, the implied
Markov regime durations (in trading days) are:

\vspace{0.2cm}

\begin{itemize}
	\item Regime~1 (low vol.): $1/(1-0{,}99753) \approx \mathbf{404}$ trading days
	\item Regime~2 (high vol.): $1/(1-0{,}99373) \approx \mathbf{160}$ trading days
\end{itemize}

Both regime \emph{durations} are longer than their respective
\emph{within-regime half-lives} --- expected, since a regime
can persist long after most of a given shock's direct effect
has already dissipated within it.

\begin{limitbox}
	$\nu_2 = 22{,}0$ ($p=0{,}096$, not significant at 5\,\%) is much
	higher than Regime~1's $\nu_1 = 7{,}07$ (strongly significant)
	--- implying near-Gaussian tails in the high-volatility
	regime, versus fat tails in the low-volatility
	regime. This is plausible (the regime split could
	itself absorb some of the extreme events that drove the
	single-regime model's fat tails, similar to the mechanism discussed
	for the HMM's Gaussian emissions,
	\S\ref{sec:regime-hmm}) but the large standard error ($16{,}85$)
	means this specific estimate should not be treated
	as precise --- likely a consequence of Regime~2 containing
	fewer observations ($\approx 28\,\%$ of 4,175 $\approx 1,183$
	days) than Regime~1. A naive $95\,\%$ Wald interval
	would give roughly $[-11, 55]$, but is not meaningful here: $\nu$
	is a bounded parameter (must exceed 2 for the variance to exist),
	and the Wald construction --- symmetric around the point
	estimate, assuming approximate normality --- respects
	neither this bound nor the typically skewed sampling
	distribution nearby. The negative lower bound should never
	be cited or interpreted.
\end{limitbox}

\subsection{Comparison with the single-regime GARCH}

\begin{table}[H]
	\centering
	\caption{MS-GARCH ($K=2$) vs single-regime GARCH comparison}
	\begin{tabular}{lrr}
		\toprule
		& Single regime (Student's $t$) & MS-GARCH ($K=2$) \\
		\midrule
		Free parameters & 5 ($\mu, \omega, \alpha, \beta, \nu$) & 10 \\
		Log-likelihood & $-2818{,}39$ & $-2806{,}64$ ($+11{,}75$) \\
		AIC & 5646.78 & \textbf{5633.27} ($-13{,}51$, favors MS-GARCH) \\
		BIC & 5678.46 & \textbf{5696.64} ($+18{,}18$, favors single regime) \\
		\bottomrule
	\end{tabular}
\end{table}

\vspace{0.3cm}

\begin{figure}[H]
	\centering
	\includegraphics[width=1\linewidth]{03_msgarch_vs_single_regime.png}
	\caption{MS-GARCH vs simple GARCH --- the smoothing the MS-GARCH provides.}
	\label{fig:msgarch-vs-single}
\end{figure}

\vspace{0.3cm}

\textbf{AIC and BIC disagree} --- a common and expected result
when comparing a regime-switching model to its single-regime
counterpart: AIC penalizes extra parameters less severely
than BIC, and the MS-GARCH adds 5 parameters relative to the
single-regime model (10 vs 5; the single-regime fit
also estimates a constant mean that the MS-GARCH
specification does not estimate). This is reported as such, a genuine
ambiguity, not resolved in favor of either model --- the honest
conclusion is that \textbf{the MS-GARCH's added complexity buys a
	real but modest fit improvement, at a parameter cost
	that BIC judges not worthwhile, while AIC judges it
	worthwhile}. Neither
criterion is "more correct" in general; this is a
known, unresolved tension in model selection between AIC
(better for prediction-oriented use) and BIC (better for
parsimony~/~true-model-identification-oriented use). The raw gap
between the two criteria is small in absolute terms at this
scale ($\approx 5,600$); a plot of the gap directly
($\Delta$AIC $=-13{,}51$, $\Delta$BIC $=+18{,}18$) rather than raw
bars is in the appendix (Figure~\ref{fig:annex-info-criteria-delta}).

\newpage

\subsection{Case study: does the model detect the COVID-19 shock (March 2020)?}

This check was carried out specifically because an initial
look at the annual regime classifications (via the Viterbi
path, the most likely regime per day) showed 2020 spending
only 36 of 250 trading days ($\approx 14\,\%$) in the high-volatility
regime --- surprisingly low compared to 2010 ($100\,\%$ of days) and
2011 ($81$--$84\,\%$), years of the euro-area debt crisis, and 2022
($80$--$84\,\%$). This was investigated rather than assumed to be
either a genuine economic finding or a model flaw.

\vspace{0.3cm}

\begin{table}[H]
	\centering
	\caption{Smoothed probability of the high-volatility regime, around the shock}
	\begin{tabular}{lr}
		\toprule
		Date & $P(\text{high-vol. regime})$ \\
		\midrule
		2020-02-18 & $9{,}2\,\%$ \\
		2020-02-25 & $56{,}8\,\%$ \\
		2020-02-26 & $91{,}5\,\%$ \\
		2020-02-28 & $99{,}4\,\%$ \\
		2020-03-11 & \textbf{$99{,}998\,\%$} \\
		2020-03-19 & $98{,}4\,\%$ \\
		2020-04-01 & $77{,}6\,\%$ \\
		2020-04-17 & $25{,}4\,\%$ \\
		2020-04-30 & $14{,}9\,\%$ \\
		\bottomrule
	\end{tabular}
\end{table}

\begin{figure}[H]
	\centering
	\includegraphics[width=1\linewidth]{06_covid_detection_zoom.png}
	\caption{Trajectory of the smoothed stress-regime probability around the COVID shock --- the table's 9 milestone dates, visualized.}
	\label{fig:covid-detection-zoom}
\end{figure}

\newpage

\textbf{In hindsight, the model identifies the shock clearly and
	quickly} --- the probability rises from $9\,\%$ to over $99\,\%$ in
about three weeks (February~18 -- March~11), tracking exactly the
actual COVID market panic, then declines over the following
six weeks.

\vspace{0.6cm}

\begin{limitbox}
	\textbf{These are smoothed probabilities, not a real-time signal.}
	The value shown for a given day uses observations
	up to the end of the sample, including the days that followed. The
	table describes \emph{when} the regime switched, not what a
	user would have known on that day. Real-time behavior is
	measured in section~\ref{sec:covid}, with a model estimated on data
	strictly prior to 2020: the predicted probability of the
	stress regime at forecast time was $0{,}3\,\%$ on February 21 and 27,
	$2{,}5\,\%$ on March 2, $49{,}3\,\%$ on March 6, and $98{,}7\,\%$
	on March 12. This is not the same model or the same quantity (predicted,
	not smoothed), so the two series are not strictly
	comparable, but a lag of about one to two weeks relative
	to the smoothed trajectory above is the order of magnitude
	relevant for an operational user.
\end{limitbox}

\vspace{0.3cm}

\textbf{Why the annual Viterbi count was misleading.} The
Viterbi path only records which regime has the
\emph{highest} probability on a given day (a binary $>50\,\%$
threshold), reducing a smooth, informative probability
trajectory to a crude yes~/~no label. The high-probability window here (February~25
-- April~9, 33 consecutive trading days above $50\,\%$) is
close to the 36-day count found --- the annual count was
not wrong, but it \textbf{discards the shape of the transition}: a
genuinely fast, extreme spike ($9\,\%\to99{,}998\,\%$ in three
weeks) looks statistically identical, in a simple daily
count, to a slower, weaker spike that crosses $50\,\%$
for a similar number of days. 
\vspace{0.2cm}
\textbf{Interpretive lesson,
	documented here rather than left implicit}: comparing regimes by
days-per-year count is too coarse a measure to judge the
relative intensity of a shock; the smoothed probability
trajectory (peak level, transition speed) is the economically
meaningful signal, and the Viterbi~/~daily-count view should be treated
as a simplification, not a substitute.

\vspace{0.3cm}

\begin{figure}[H]
	\centering
	\includegraphics[width=1\linewidth]{05_annual_regime_share.png}
	\caption{Annual stress-regime share, under three definitions (Viterbi, smoothed $>50\,\%$, average probability) --- resolves the 39--100\,\% gap that the initial screening left open.}
	\label{fig:annual-regime-share}
\end{figure}

\subsection{Economic interpretation}

\begin{econbox}
	\begin{itemize}
		\item EUR/USD alternates between two distinct volatility
		regimes: a \textbf{calm regime} (realized daily
		volatility $\approx 0{,}44$ on days when it dominates, occurring
		$\approx 72\,\%$ of the time, lasting on average about
		404 trading days, or $\approx$~1.6 years, once entered) and
		a \textbf{stress regime} ($\approx 0{,}70$ on days when it
		dominates --- about 1.6 times more volatile --- occurring
		$\approx 28\,\%$ of the time, lasting on average about
		160 trading days, or $\approx$~7--8 months). The implied
		long-run $\sigma$ from the parameters ($0{,}21$ vs $0{,}73$)
		overstates this gap and is imprecise.
		\item In the point estimates, shock intensity within a
		stress regime fades \textbf{faster} (half-life
		$\approx 64$ days) than within a calm regime
		($\approx 237$ days) --- stress would be acute but
		comparatively brief at the shock level, even though the
		\emph{regime itself} (the Markov state) can still take
		months to fully resolve. This difference is not
		statistically established and should be read as suggestive.
		\item In hindsight, the model places the genuine, sudden
		crisis (COVID, March 2020) in about three weeks, moving
		from near-certain "calm" to near-certain
		"stress". This is a retrospective description of the
		moment the regime switched; as an operational
		early-warning signal, it lagged by about one to two
		weeks in the out-of-sample test (section~\ref{sec:covid}).
		\item Compared with a single, undifferentiated GARCH
		(section~\ref{sec:garch}), this model offers a genuinely
		different, regime-aware description of volatility: it
		explicitly answers "which of the two regimes are we in?",
		rather than averaging the dynamics of both regimes into a single
		fixed set of parameters. The formal test in
		section~\ref{sec:garch-stability} did \textbf{not} establish that the
		single-regime model's parameters are statistically
		unstable over time (likelihood-ratio test,
		$p=0{,}277$): the MS-GARCH's motivation rests more solidly
		on being the theoretically appropriate tool for
		regime-dependent dynamics and on its AIC advantage than on
		proven instability of the simpler model --- at the acknowledged
		cost of a heavier, harder-to-estimate model, whose
		advantage is judged worthwhile by AIC and not by BIC.
	\end{itemize}
\end{econbox}

\subsection{Known limitations}

\begin{limitbox}
	\begin{itemize}
		\item MSGARCH's likelihood is exact for the (more
		restricted) specification of \citet{haasmittnikpaolella2004} that it
		implements --- this is not an approximation of that model. What
		remains genuinely intractable is the fully general,
		path-dependent MS-GARCH likelihood; MSGARCH avoids
		this intractability by using a different, deliberately more
		restricted model, not by approximating the
		general model.
		\item $\nu_2$ (regime~2's tail parameter) is imprecisely
		estimated; should not be over-interpreted alone.
		\item The persistence difference between the two regimes is not
		tested; the half-lives are point estimates,
		sensitive to $\alpha+\beta$ being close to 1.
		\item The implied long-run $\sigma$ figures from the
		parameters ($0{,}206$ and $0{,}730$) rely on $\omega$, not
		significant in either regime, and on a denominator
		$1-\alpha-\beta$ close to zero: they are imprecise, and their
		ratio ($3{,}5$) is well above the ratio of realized
		volatilities on days when each regime dominates ($1{,}6$).
		\item AIC and BIC disagree on whether the added
		complexity is justified --- this project does not settle
		between the two criteria, reporting the disagreement rather than
		picking a side to force a clean narrative. The two models also
		differ in their mean specification (a constant mean
		estimated in the single-regime fit, none here).
		\item The regime probabilities analyzed above are smoothed
		(\emph{ex post}). A real-time user would rely on
		filtered or predicted probabilities, which react more
		slowly.
		\item \textbf{The Viterbi path used above is not produced
			by any script currently in the project.}
		\texttt{ms\_garch.R}, as written, only exports the
		smoothed regime probabilities
		(\texttt{ms\_garch\_regime\_probs.csv}) from
		\texttt{State()}'s output; the Viterbi classification behind the
		opening paragraph and the annual day counts was generated
		in an interactive session and not captured in a saved
		script. Anyone reproducing this analysis from the
		repository alone would need to add a Viterbi path
		extraction from \texttt{State()}'s output themselves --- not yet
		done.
		\item $K=2$ has not been compared to $K=3$ at this stage (see
		the justification for starting at 2 above); a
		natural extension, not yet undertaken.
		\item This step depends on R and the MSGARCH package --- a
		deliberate, documented exception to the rest of the
		project's Python pipeline, meaning it cannot be re-run
		purely from the Python codebase without R installed
		alongside it.
	\end{itemize}
\end{limitbox}

\vspace{0.4cm}

\begin{limitbox}
	\textbf{Correction made to the present write-up:} an earlier
	version of this document flagged, at this point, the absence of a file for the
	within-regime half-life~/~Markov regime-duration
	distinction --- a search error in the archive, corrected after a second
	exhaustive pass. The file exists (figure below),
	from an earlier production batch predating the
	post-audit corrections of this document.
\end{limitbox}

\vspace{0.4cm}

\begin{figure}[H]
	\centering
	\includegraphics[width=0.8\linewidth]{03_halflife_vs_regime_duration.png}
	\caption{Within-regime half-life (dissipation of GARCH shocks) vs regime duration in the Markov sense --- this section's subtlest distinction, accompanied by the figure rather than left to text alone.}
	\label{fig:halflife-vs-regime-duration}
\end{figure}

\newpage

% ═════════════════════════════════════════════════════════════════
\section{Operational price range}
\label{sec:price-range}
% ═════════════════════════════════════════════════════════════════

\subsection{Objective}

Produce an actual computed price range
$[P_{\min}, P_{\max}]$ for EUR/USD, using the fitted
MS-GARCH($K=2$) model (section~\ref{sec:msgarch}) --- not the
illustrative made-up figures used earlier in the project to explain
the \emph{concept} of a regime-aware range
(sections~\ref{sec:magnitude} and~\ref{sec:regime-hmm}), but a value
actually computed from the fitted model and the current real
data.




\subsection{Results}
\label{sec:price-range-results}

As of 2026-09-11 (EUR/USD $=1{,}1604$):

\vspace{0.3cm}

\begin{center}
	\begin{tabular}{|p{6.7cm}|p{6.7cm}|}
		\hline
		Current regime probabilities & low vol. $=99{,}4\,\%$, high vol. $=0{,}6\,\%$ \\
		\hline
		Implied forecast volatility from the model (regime-weighted) & 0.3055 \\
		\hline
		\textbf{95\,\% price range} & \textbf{[1.1533~; 1.1675]} \\
		\hline
		Width & 0.0142 ($1{,}22\,\%$ of the current price) \\
		\hline
	\end{tabular}
\end{center}

\vspace{0.3cm}

\subsection{Internal consistency check}
\label{sec:price-range-internal-check}

The interval corresponds to about $\pm 2{,}0$ times the
forecast volatility ($-2{,}011\sigma$ / $+1{,}984\sigma$). This was checked against
the theoretical quantile of a standardized Student's~$t$
distribution with $\nu \approx 7{,}07$ (Regime~1's tail
parameter, which dominates the mixture given its $99{,}4\,\%$
weight): $-1{,}998$. The close match confirms that
\texttt{Risk()} genuinely uses the model's Student's~$t$
predictive distribution, without silently falling back to a
Normal approximation --- a genuine check, not an
assumption taken for granted.

At the $95\,\%$ level, this check has limited power to
distinguish a Student's~$t$ from a Normal: the two quantiles
differ by only $2\,\%$ ($1{,}998\sigma$ vs $1{,}960\sigma$). The
$80\,\%$ level of the \emph{fan chart} (\S\ref{sec:price-range-fanchart})
discriminates better: the mixture's $80\,\%$ bounds are at
$-1{,}212\sigma$ and $+1{,}185\sigma$, about $1{,}20\sigma$ on
average, versus $1{,}197\sigma$ for a standardized
Student's~$t$ ($\nu \approx 7{,}07$) and $1{,}282\sigma$ for a Normal.

\newpage

\subsection{Rigorous result vs genuinely regime-conditional forecast}

An earlier version of this document compared the operational
range to an \emph{approximation} of what each regime would
imply (using each regime's long-run volatility level and the
mixture's quantile multiplier, not a genuine one-step
forecast). This has since been replaced: section~\ref{sec:msgarch}
was corrected to recognize that MSGARCH's
regime-switching option implements
\citet{haasmittnikpaolella2004}, not \citet{gray1996} --- which
means the conditional-variance recursion specific to each
regime has \textbf{no} path-dependency problem and can
be computed directly:

\vspace{0.3cm}

\begin{keyeq}
	\begin{equation}
		\sigma_k^2(t) = \omega_k + \alpha_k\, r(t-1)^2 + \beta_k\, 	\sigma_k^2(t-1),
	\end{equation}
\end{keyeq}

run independently for each regime~$k$ on the same observed
return series. This gives a genuine exact one-step forecast
per regime --- not an approximation.

\textbf{Validation of the manual recursion.} Before trusting
this recursion, it was validated against the already-validated
\texttt{Risk()} mixture result: the probability-weighted average of the
two regimes' forecast variances,

\vspace{0.3cm}

\begin{center}
	
	$\sqrt{w_1\sigma_1^2(t+1) + w_2\sigma_2^2(t+1)}$
	
\end{center}

\vspace{0.3cm}
gives
\textbf{0.3055} --- exactly matching the mixture
volatility forecast reported by \texttt{Risk()} (\textbf{0.3055}).
This independent match confirms the manual recursion
is correct, rather than relying on the formula simply
being "plausible."

\vspace{0.3cm}

\begin{table}[H]
	\centering
	\caption{Genuinely regime-conditional 95\,\% ranges}
	\begin{tabular}{lrrr}
		\toprule
		& Forecast volatility $\sigma_k(t+1)$ & 95\,\% range & Width \\
		\midrule
		Regime~1 (calm) & 0.3033 & [1.1534~; 1.1675] & $1{,}21\,\%$ \\
		Regime~2 (stress) & 0.5769 & [1.1472~; 1.1737] & $2{,}28\,\%$ \\
		\textbf{Mixture (operational)} & 0.3055 & \textbf{[1.1533~; 1.1675]} & \textbf{$1{,}22\,\%$} \\
		\bottomrule
	\end{tabular}
\end{table}

\vspace{0.3cm}

The stress regime's range is genuinely (not approximately)
\textbf{1.88 times wider} than the calm regime's --- a
direct consequence of Regime~2's own one-step volatility
being almost double Regime~1's ($0{,}577$ vs
$0{,}303$, a ratio of $1{,}90$). The tail shape plays only a
minor role, and in the opposite direction: Regime~2's tail is
\emph{lighter} ($\nu_2 \approx 22$, closer to Normal than
expected, section~\ref{sec:msgarch}) than Regime~1's
($\nu_1 \approx 7$), which very slightly reduces the ratio at the
$95\,\%$ level ($1{,}88$ vs a volatility ratio of $1{,}90$).

\vspace{0.3cm}

\begin{limitbox}
	The $95\,\%$ interval reported by \texttt{Risk()} for the mixture
	($-0{,}6143$ / $+0{,}6061$) is very slightly asymmetric, even
	though the model has no mean term at all (\texttt{MSGARCH}
	assumes zero mean everywhere) and each individual per-regime
	Student's~$t$ distribution is exactly symmetric
	around zero. A mixture of two symmetric distributions centered
	at zero is itself theoretically symmetric, and its exact
	$95\,\%$ quantiles, computed from the published parameters and
	current regime weights, are $\pm 0{,}611$. The reported
	bounds deviate from this value by less than $1\,\%$, and the
	deviation is systematic rather than random: at every
	\emph{fan chart} confidence level, the midpoint of the lower
	and upper bounds is the same, $-0{,}0041$, and each bound reported by
	\texttt{Risk()} sits on a regular numerical grid with step
	$\approx 0{,}0068$ percentage points of return. The asymmetry is
	therefore a grid effect of \texttt{Risk()}'s numerical computation --- not
	Monte Carlo noise, and not a property of the model. It is
	negligible for the price range ($\approx 0{,}00005$ in price).
\end{limitbox}


\vspace{0.3cm}

This is no longer a cosmetic improvement over the earlier
approximation: the genuine regime-conditional ranges are
computed with the same rigor as the operational mixture range,
using each regime's own parameters and its own
tail shape --- not borrowed multipliers or long-run
levels. This makes the regime comparison itself
a legitimate, citable result, not a mere illustrative aside.

\newpage

\subsection{What "next trading day" really means --- a clarification that avoids a false failure}

\begin{limitbox}
	\texttt{nahead = 1} means one step forward \textbf{within the
		series of trading days used to fit the model}, not literally
	"24 hours from now." The last observation used was
	\textbf{Friday 2026-09-11}; reference FX series like
	FRED's \texttt{DEXUSEU} have no observation on
	weekends. The range above is therefore a forecast for the
	\textbf{next available trading-day observation} --- here,
	\textbf{Monday 2026-09-14} --- not Saturday.
	
	\textbf{An additional, more substantial caveat}: the actual FX
	market trades continuously from Sunday evening to Friday
	evening (New York time), unlike the single daily
	\texttt{DEXUSEU} reference rate (a New York noon rate) on
	which the model was calibrated. A whole weekend of news
	(central bank comments, geopolitical events) can occur between
	Friday's close and Monday's reopening, which the model has
	\textbf{never observed as a distinct event} --- it was
	calibrated purely on day-to-day changes in this
	reference rate, with no explicit "weekend gap"
	component. In practice: a Monday's actual move
	could fall outside the $95\,\%$ range more often than an
	"average" weekday, purely because weekends can carry
	more unobserved information than a simple weekday
	gap. This is a known limitation of daily-close
	GARCH-family models applied to FX, not a flaw specific
	to this implementation --- but it should be stated
	explicitly rather than left for a reader to discover by
	observing the model "fail" on a Monday that, in fact,
	behaved exactly as a model blind to weekend news
	would be expected to.
\end{limitbox}

\newpage

\subsection{Economic interpretation}

\begin{econbox}[Illustrative example]
	\textbf{A wide $1{,}22\,\%$ range over a single day is a
		real, actionable figure, not an abstraction.} For a
	firm converting a 1,000,000~EUR invoice into USD, the
	range $[1{,}1533~;~1{,}1675]$ implies the dollar proceeds
	could plausibly fall anywhere between about
	\$1,153,300 and \$1,167,500 --- a gap of more than \$14,000 due to
	exchange-rate uncertainty alone over a single day, before any
	other business risk. This is the kind of figure a treasury or
	finance function would use to size a hedge, not
	a mere academic quantity.
	
	\textbf{The width of the range is not fixed --- it
		"breathes" with the regime}, which is the whole
	operational point of this project. As of this date, the model assigns
	$99{,}4\,\%$ probability to the calm regime, so the
	operational range is close to its calm-regime width.
	Using the exact regime-conditional ranges above, the same
	1,000,000~EUR conversion directly illustrates the
	difference:
	\begin{itemize}
		\item \textbf{If calm (certain):} $[1{,}1534~;~1{,}1675]$ $\to$
		proceeds between \$1,153,400 and \$1,167,500 --- a gap of
		\$14,100, almost identical to the operational range, since
		the model is currently near-certain the market is calm.
		\item \textbf{If stress (certain):} $[1{,}1472~;~1{,}1737]$ $\to$
		proceeds between \$1,147,200 and \$1,173,700 --- a
		\textbf{gap of \$26,500}, almost double. A treasury relying
		on a single, regime-blind volatility estimate calibrated
		during calm periods would be under-hedged by roughly this
		difference at the very moment the market actually shifts into
		stress --- exactly when the hedge matters most.
	\end{itemize}
	
	\textbf{What the range does NOT say:} nothing about direction. A
	$95\,\%$ range of $[1{,}1533~;~1{,}1675]$ is symmetric
	information about \emph{how far} EUR/USD might move, not
	\emph{which way} --- consistent with the project's earlier
	robust finding that direction is not predictable at this
	horizon (sections~\ref{sec:baseline} and~\ref{sec:controls}). Anyone
	reading this range as "the model expects EUR/USD to end up
	somewhere here, probably in the middle" is
	over-interpreting it: the model has no opinion on where, within the
	range, the price will land, only that $95\,\%$ of the time, it
	expects the actual move to fall inside. This holds
	in both regimes: a wider stress-regime range does not mean the
	model expects a crash more than a rally --- only a
	wider move in one direction or the other.
\end{econbox}

\begin{figure}[H]
	\centering
	\includegraphics[width=1\linewidth]{05_economic_impact_eur1m.png}
	\caption{Economic impact of a 1,000,000~EUR conversion --- \$14,100 (calm) vs \$26,500 (stress).}
	\label{fig:economic-impact}
\end{figure}

\subsection{Ex-post comparison against actual market data \\(2026-09-14)}

Independent source (Pound Sterling Live's daily EUR/USD
history), for Monday 2026-09-14 (the next trading day after the
2026-09-11 forecast):

\begin{center}
	\begin{tabular}{lr}
		\toprule
		& Value \\
		\midrule
		Open & 1.1597 \\
		Close & 1.1549 \\
		High & 1.1601 \\
		Low & 1.1523 \\
		\bottomrule
	\end{tabular}
\end{center}

\vspace{0.3cm}

Forecast range: \textbf{[1.1533~; 1.1675]}.


Open, Close, and High all fall inside the
forecast range. The \textbf{intraday Low (1.1523) falls
	just outside the lower bound}, by 0.0010 ($\approx 0{,}086\,\%$
of the reference price --- about ten pips).

\newpage

\begin{limitbox}
	\textbf{The most relevant comparison is against the Close (1.1549),
		not the intraday Low.} The GARCH was calibrated on the
	daily FRED reference rate's day-to-day log returns --- it
	was never built to predict intraday extremes
	(the high/low spread within a single day), only
	the day-to-day move size. Judged on that quantity, the
	Close falls comfortably within the range: consistent with
	the model, though only indicatively (see the sourcing caveat
	below). The Low breaching the band is a data point
	on a different question (intraday range) that this
	model does not address --- worth noting, not a failure of
	what was actually forecast.
	
	\textbf{Sourcing caveat:} this project's own pipeline
	recorded the 2026-09-11 reference price as 1.1604
	(FRED's \texttt{DEXUSEU}), while this independent source
	reports a Close of 1.1599 for the same day --- a gap
	of about $0{,}04\,\%$. The two are not the same
	quantity: \texttt{DEXUSEU} is the Federal Reserve's New York
	noon reference rate, not an end-of-day close, so the
	model forecasts the next noon reference rate while
	Pound Sterling Live reports the end-of-day close. The gap is
	small on a calm day but can be much wider during
	volatile periods (averaging $0{,}22\,\%$ between February 24 and
	April 30, 2020, up to $1{,}14\,\%$ --- section~\ref{sec:covid}), so
	this comparison is indicative, not exact.
	
	\textbf{Critical caveat: this is an observation, not a
		validation.} A single day proves essentially nothing about
	the model's calibration. A correctly calibrated $95\,\%$
	interval is \emph{supposed} to be breached roughly 1 day in
	20 --- a single result inside the range (or even outside) carries no
	statistical weight on its own. This comparison is recorded as an
	honest, concrete illustration of the range's performance
	on one occasion, not as proof that the model is well or
	poorly calibrated. A genuine calibration check would
	require accumulating many such observations over time --- not
	attempted here (see section~\ref{sec:rolling-test} for this
	accumulation).
\end{limitbox}

\newpage

\subsection{Fan chart: multiple confidence bands}
\label{sec:price-range-fanchart}

Extending the single $95\,\%$ range into a \emph{fan chart} of several nested confidence bands
computed in a single \texttt{Risk()} call over the full alpha
vector, converted into price bands.

\vspace{0.3cm}

\begin{table}[H]
	\centering
	\caption{Band widths, three panels}
	\begin{tabular}{|p{2.5cm}|p{2.5cm}|p{2.5cm}|p{2.5cm}|p{2.5cm}|}
		\hline
		Confidence level & Mixture width & Calm width & Stress width & Stress/calm ratio \\
		\hline
		20\,\% & 0.136\,\% & 0.135\,\% & 0.282\,\% & 2.09$\times$ \\
		\hline
		50\,\% & 0.366\,\% & 0.365\,\% & 0.754\,\% & 2.07$\times$ \\
		\hline
		80\,\% & 0.732\,\% & 0.726\,\% & 1.454\,\% & 2.00$\times$ \\
		\hline
		95\,\% & 1.220\,\% & 1.212\,\% & 2.282\,\% & 1.88$\times$ \\
		\hline
	\end{tabular}
\end{table}

\vspace{0.3cm}

The bands are correctly nested at every level
($20\,\% \subset 50\,\% \subset 80\,\% \subset 95\,\%$) and widen
faster than linearly toward the tails in all three panels,
as a Normal distribution would too. What distinguishes the
model is the degree: the $95\,\%$ band is about $9{,}0$ times
wider than the $20\,\%$ band in the mixture and calm
panels ($8{,}1$ times in the stress panel), versus $7{,}7$ times for a
Normal distribution --- the signature of the fatter tails of the
Student's~$t$ distribution underlying the model.

\vspace{0.4cm}

\begin{figure}[H]
	\centering
	\includegraphics[width=1\linewidth]{01_fan_chart_three_panels.png}
	\caption{Three-panel \emph{fan chart} (operational mixture, certain calm regime, certain stress regime).}
	\label{fig:fan-chart}
\end{figure}

\begin{figure}[H]
	\centering
	\includegraphics[width=1\linewidth]{02_width_ratio_stress_calm.png}
	\caption{Stress~/~calm width ratio by confidence level --- from $2{,}09\times$ to $1{,}88\times$, not constant.}
	\label{fig:width-ratio-stress-calm}
\end{figure}

\vspace{0.3cm}

\subsection{Interpretation --- the calm regime, and, especially, the stress regime.}

\begin{econbox}[Calm Regime]
	The calm regime is nearly identical to the operational mixture
	at every level (expected, since the mixture is $99{,}4\,\%$
	calm-weighted as of this date) --- confirming that the mixture
	range is, for all practical purposes, "the calm-regime range" at this
	moment.
\end{econbox} 

\vspace{0.2cm}

\begin{econbox}[Stress Regime: what this chart's panel shows]
	At every confidence level, the
	stress-regime band is \textbf{about twice as wide} as
	the calm-regime band --- but this ratio is not constant, and
	how it changes is itself informative. It is
	\textbf{$2{,}09\times$, $2{,}07\times$, and $2{,}00\times$} at the
	$20\,\%$, $50\,\%$, and $80\,\%$ levels, about $5$--$10\,\%$
	\emph{above} the ratio of the two regimes' one-step volatility
	forecasts ($0{,}577/0{,}303 \approx 1{,}90$), and \textbf{$1{,}88\times$} at
	the $95\,\%$ level, essentially equal to it. If the two regimes
	had the same distribution shape, the ratio would be constant and equal
	to $1{,}90$; the deviations measure the effect of the different
	shapes documented in section~\ref{sec:msgarch}. At a given variance, a
	fat-tailed distribution places more mass near the center,
	so the calm regime's central band is narrower than
	a near-Normal distribution would give: at the $20\,\%$
	level, the half-width is $0{,}223\sigma$ for the calm regime ($\nu_1
	\approx 7$) versus $0{,}245\sigma$ for the stress regime ($\nu_2
	\approx 22$; a Normal gives $0{,}253\sigma$). This inflates the
	stress/calm ratio at central levels. In the extreme tail, the two
	shapes give nearly the same multiple of $\sigma$ ($1{,}998\sigma$
	versus $1{,}977\sigma$ at $95\,\%$), and the ratio reverts to the
	volatility ratio. The \emph{absolute gap} between the two bands, on the
	other hand, keeps growing with the confidence level....
\end{econbox}


\vspace{0.2cm}

\begin{econbox}[What this means for someone reading the chart, not just the figures]
	The stress panel is not simply "the
	calm panel stretched by a constant factor" --- it has a
	genuinely different shape, relatively narrower at its most
	extreme edge than at its center, because the two regimes have
	different distribution shapes. A reader looking only at the
	$95\,\%$ width would actually \emph{underestimate} slightly
	how much riskier the stress regime's central, day-to-day
	experience is compared with the calm regime's (the $20\,\%$ and $50\,\%$
	ratios, $\approx 2{,}1\times$, are larger than the headline $95\,\%$
	ratio, $\approx 1{,}88\times$) --- something to bear in mind before using a simple
	"times two" rule of thumb to reason about regime risk at
	any horizon.
\end{econbox}


\subsection{Status of project objective (b)}

With this step, objective (b) --- stated from the very start of the project,
\S\ref{sec:intro} --- is now concretely delivered: an
actual, computed, regime-aware price range, anchored in a
validated GARCH specification (section~\ref{sec:garch}), an honestly
labeled regime-detection layer
(sections~\ref{sec:regime-hmm} and~\ref{sec:msgarch}), and a
documented, verified forecasting method (this section) --- rather than the
illustrative example used earlier in the project purely to explain the
concept.

\newpage

% ═════════════════════════════════════════════════════════════════
\section{Rolling chain test (September 11--17, 2026)}
\label{sec:rolling-test}
% ═════════════════════════════════════════════════════════════════

\subsection{Objective and design}

Chain together 4 sequential one-step forecasts, each using only
the information available up to the previous trading day:

\begin{table}[H]
	\centering
	\caption{Rolling test design}
	\begin{tabular}{lll}
		\toprule
		Forecast & History up to & Target \\
		\midrule
		1 & 2026-09-11 & 2026-09-14 \\
		2 & 2026-09-14 & 2026-09-15 \\
		3 & 2026-09-15 & 2026-09-16 \\
		4 & 2026-09-16 & 2026-09-17 \\
		\bottomrule
	\end{tabular}
\end{table}

\vspace{0.3cm}

\textbf{The parameters are frozen} (the MS-GARCH fit from
\texttt{ms\_garch.R}, unchanged) --- each forecast differs only
because the return history feeding \texttt{State()} and
\texttt{Risk()} grows by one actually observed day each time. This
isolates exactly one question: how would the model, calibrated
once on 2026-09-11, have tracked reality over the following
week without being retrospectively re-fitted.

\subsection{Data source workaround}

A fresh FRED (\texttt{DEXUSEU}) pull turned out to still stop
at 2026-09-11 --- the same end date as the project's
original pull --- either a genuine publication lag or a
query issue; unresolved at the time of writing. As a
documented workaround, the existing FRED history was extended with
independently sourced Close values (Pound Sterling Live) for
2026-09-14 through 2026-09-17 --- the same source and the same gap already noted
for 2026-09-11 (section~\ref{sec:price-range}) in the price
range document.

\subsection{Results (close to close)}

\begin{table}[H]
	\centering
	\small
	\caption{Rolling test results, close to close}
	\begin{tabular}{|p{1.5cm}|p{1.5cm}|p{1.5cm}|p{1.5cm}|p{1.5cm}|p{1.5cm}|p{1.5cm}|p{1.5cm}|}
		\hline
		Forecast & Threshold price & Range & Width & Actual (Close) & Result & $P$(calm) & $P$(stress) \\
		\hline
		1 ($\to$09-14) & 1.1604 & [1.1533~; 1.1675] & 1.22\,\% & 1.1549 & \textbf{Inside} & 99.4\,\% & 0.6\,\% \\
		\hline
		2 ($\to$09-15) & 1.1549 & [1.1477~; 1.1620] & 1.23\,\% & 1.1542 & \textbf{Inside} & 98.9\,\% & 1.1\,\% \\
		\hline
		3 ($\to$09-16) & 1.1542 & [1.1471~; 1.1613] & 1.23\,\% & 1.1464 & \textbf{Outside} & 99.2\,\% & 0.8\,\% \\
		\hline
		4 ($\to$09-17) & 1.1464 & [1.1390~; 1.1538] & 1.29\,\% & 1.1459 & \textbf{Inside} & 97.0\,\% & 3.0\,\% \\
		\hline
	\end{tabular}
\end{table}

\textbf{Raw Close-based result: 3 of 4 inside the
	nominal $95\,\%$ interval.}

\vspace{0.7cm}
\begin{figure}[H]
	\centering
	\includegraphics[width=1\linewidth]{02_range_vs_actual_chain.png}
	\caption{Forecast range vs realized price, chained across the 4 forecasts.}
	\label{fig:range-vs-actual-chain}
\end{figure}

\textbf{Adding Open, High, Low --- a fuller, slightly less
	favorable picture.} The model is calibrated on the
day-to-day changes of a single daily reference rate, so the
daily Close is the closest primary comparison. But
checking the full daily range against each forecast
band reveals a nuance invisible from the Close alone:

\vspace{0.3cm}	

\begin{table}[H]
	\centering
	\small
	\caption{Open, High, Low, Close against range}
	\begin{tabular}{lllll}
		\toprule
		Target & Range & Open & High & Low \\
		\midrule
		09-14 & [1.1533~; 1.1675] & 1.1597 (inside) & 1.1601 (inside) & \textbf{1.1523 (OUTSIDE)} \\
		09-15 & [1.1477~; 1.1620] & 1.1549 (inside) & 1.1553 (inside) & 1.1527 (inside) \\
		09-16 & [1.1471~; 1.1613] & 1.1542 (inside) & 1.1557 (inside) & \textbf{1.1461 (OUTSIDE)} \\
		09-17 & [1.1390~; 1.1538] & 1.1464 (inside) & 1.1473 (inside) & 1.1457 (inside) \\
		\bottomrule
	\end{tabular}
\end{table}

\vspace{0.3cm}

\textbf{On the Close alone: 1 breach out of 4. Including the
	intraday Low: 2 breaches out of 4} (09-14 and 09-16).
The 09-14 Low (1.1523) dipped below the forecast's lower
bound even though that day's Close came back inside by the
end of the session. This does not change the model's validity (it was never
built to predict intraday extremes), but it is
a fuller picture than reporting only the Close-based
result.



\subsection{Interpretation --- corrected before writing, not after}


\begin{limitbox}
	\textbf{$n=4$ proves nothing about calibration --- stated
		plainly, not softened.} 3/4 (or 2/4, including the intraday
	Lows) inside a $95\,\%$ interval neither validates
	nor invalidates the model's calibration. A meaningful empirical
	coverage rate requires at minimum several dozen observations
	before it carries statistical information. \textbf{This result is
		a demonstration that the forecast chain works correctly
		end to end, not proof of good (or bad) calibration.} Any
	claim that "the model captures uncertainty well"
	based on this alone would be an unsupported generalization
	from a tiny sample.
	
	For scale: even if the $95\,\%$ interval were perfectly
	calibrated and breaches were independent, the probability
	of observing at least one breach in 4 days is
	$1-0{,}95^4 \approx 18{,}5\,\%$ --- a single breach in 4 days
	is an ordinary result.
	
	
\end{limitbox}

\newpage

\subsubsection{The Close-based miss (forecast~3, $\to$09-16): two distinct explanations, correctly kept separate}

The breach was marginal: actual $1{,}1464$ against the lower
bound $1{,}1471$, a gap of $0{,}0007$ ($\approx 0{,}064\,\%$ of the
price, about 7 pips) --- right at the edge, not a spectacular
failure.

Two distinct hypotheses were considered to determine whether this
represents a genuine model shortfall:

\vspace{0.2cm}

\textbf{(a) Direct measurement noise on the "actual" value
	itself.} The 09-16 forecast is compared against the close reported
by Pound Sterling Live, whereas the model was calibrated on
FRED data throughout its history. Since the FRED value for
09-16 is unavailable, \textbf{it cannot be ruled out} that the
FRED value for that day would have fallen inside the range. Over
the calm January-February 2020 days (section~\ref{sec:covid}), the
absolute gap between the two sources averaged $0{,}069\,\%$
(std.\ dev.\ $0{,}088\,\%$), and 17 of those 35 days showed a gap
larger than this breach ($0{,}064\,\%$); the one-day figure
cited for 09-11 ($\approx 0{,}04\,\%$) is a single observation.
Moreover, the FOMC statement on September 16 was released at 2:00~PM
(Eastern time) --- after a noon reference rate and before an
end-of-day close --- so a noon value that day would not
contain the market reaction, unlike the close.
\textbf{This remains genuinely unresolved with the available
	data}, but the breach falls well within the range
of gaps normally observed between the two quantities: it does not,
on its own, point to a model failure.

\vspace{0.2cm}

\textbf{(b) Contamination of the GARCH recursion via the
	one-off source switch (09-11$\to$09-14) --- a distinct,
	quantifiable channel.} The single switch from FRED
(09-11 $=1{,}1604$) to Pound Sterling Live (09-14 $=1{,}1549$) creates a
"phantom" distortion in that day's computed return:

\begin{center}
	\small
	\begin{tabular}{|p{7.8cm}|p{6.5cm}|}
		\hline
		Actually computed return (FRED $\to$ PSL) & $\ln(1{,}1549/1{,}1604) = -0{,}4751\,\%$ \\
		\hline
		True single-source return (PSL 09-11 close & \\
		$=1{,}1599$ $\to$ PSL 09-14) & $\ln(1{,}1549/1{,}1599) = -0{,}4320\,\%$ \\
		\hline
		Phantom distortion & $-0{,}0431$ percentage points (larger magnitude) \\
		\hline
	\end{tabular}
\end{center}

\vspace{0.3cm}	

\begin{figure}[H]
	\centering
	\includegraphics[width=1\linewidth]{canal_garch_phantom_distortion.png}
	\caption{Flow diagram of channel (b) --- two computation paths compared side by side, $-0{,}4751\,\%$ vs $-0{,}4320\,\%$.}
	\label{fig:garch-phantom-distortion}
\end{figure}

Because the GARCH's variance recursion uses squared past
returns 

\vspace{0.2cm}

\begin{center}
	($\sigma^2(t) = \omega + \alpha\, r^2(t-1) + \beta\,
	\sigma^2(t-1)$)
\end{center}

\vspace{0.2cm}

 a \textbf{larger}-magnitude return injects
\textbf{more}, not less, variance into the forecast for each
following day --- meaning this channel, if it has an effect, would have made
forecast~3's range \textbf{wider} than an entirely
single-source computation would have given. \textbf{This
	channel does not explain the breach; if it has an effect, it
	suggests the "true" single-source model might have missed by
	slightly more, not less.}
	
\vspace{0.2cm}	

\textbf{Conclusion:} channel~(a) remains a genuine, unresolved
source of doubt; channel~(b) was quantitatively checked and does
not "exonerate" the model --- it points the other way.

\newpage

\subsubsection{Regime responsiveness: two competing hypotheses, correctly left undecided}
\label{sec:rolling-regime-reactivity}

Regime probabilities across the chain: $P(\text{stress}) = 0{,}6\,\%
\to 1{,}1\,\% \to 0{,}8\,\% \to \mathbf{3{,}0\,\%}$, while the
range width barely moved ($1{,}22\,\% \to 1{,}23\,\% \to
1{,}23\,\% \to 1{,}29\,\%$; $+5{,}6\,\%$ between the first and
last). The model's risk sentiment rose about 5-fold
($0{,}56\,\% \to 3{,}01\,\%$) over the week but stayed
anchored above $97\,\%$ "calm" throughout, including
immediately after the miss.

\vspace{0.3cm}

\begin{figure}[H]
	\centering
	\includegraphics[width=1\linewidth]{02_stress_probability_vs_width.png}
	\caption{$P(\text{stress})$ and range width, indexed to base~100 --- the actual gap ($\times 5{,}3$ vs only $+5{,}6\,\%$) stays in the subtitle to avoid any misleading effect from the indexing.}
	\label{fig:stress-prob-vs-width}
\end{figure}

\vspace{0.4cm}

\begin{limitbox}
	\texttt{Risk()} returns its bounds on a numerical grid with step
	$\approx 0{,}0068$ percentage points of return
	(\S\ref{sec:price-range-internal-check}), about $0{,}56\,\%$ of a
	$95\,\%$ width. The difference between the second and third
	widths ($1{,}234\,\%$ vs $1{,}227\,\%$) is about one grid
	step and should not be interpreted; the rise to $1{,}29\,\%$ is
	about ten steps and constitutes a genuine change.
\end{limitbox}

\newpage	

Two explanations are both consistent with this single
observation, and \textbf{cannot be distinguished with only
	4 data points}:
\begin{enumerate}
	\item \textbf{Behavior expected by design}: frozen
	parameters mean the forecast only shifts because
	the information set changes marginally each day --- no
	violent reaction is expected or desired from a single new
	observation.
	\item \textbf{A structural slowness to react to emerging
		stress}, which would only show up as a genuine problem
	(systematic under-coverage) over a longer,
	genuinely volatile period --- not detectable over a relatively
	mild week.
\end{enumerate}

%\textbf{Proposed next step to decide between the two:}
%re-run the same rolling chain with frozen parameters over a	historical window containing a real, sustained shock (e.g.\ February--April 2020) and check whether the empirical coverage rate collapses. This has since been done (section~\ref{sec:covid}), with a model estimated on data prior to 2020 only: coverage of $77/84 = 91{,}7\,\%$ on the FRED rate (one-sided binomial test $p=0{,}127$), so no collapse, but the test has limited power. Its breaches occurred both before the stress probability rose ($P(\text{stress})$ below $50\,\%$) and after it exceeded $95\,\%$, so a lag in the early days of a shock cannot be ruled out. The question is narrowed, not settled.

\newpage

\subsection{Economic interpretation}

\begin{econbox}[A concrete illustration, not just an abstract test.]
	
	Take forecast~3 ($\to$09-16), the one that missed: for a
	1,000,000~EUR conversion decided at the 09-15 close, the
	model's $95\,\%$ range implied proceeds between
	\$1,147,100 and \$1,161,300. The actual outcome (\$1,146,400) fell
	about \$730 below even the pessimistic end of that
	range --- on a \$1.15 million transaction, a shortfall
	equivalent to about $0{,}06\,\%$ of the notional. Framed this way,
	even this chain's only "miss" was a
	\textbf{near-miss in dollar terms}, not a scenario that would have
	caused significant, unplanned damage to a firm.
\end{econbox}

\vspace{0.4cm}	
\begin{figure}[H]
	\centering
	\includegraphics[width=1\linewidth]{05_economic_near_miss.png}
	\caption{Forecast~3's miss in dollar terms --- \$1,146,400 realized vs \$1,147,100 at the pessimistic end of the range, a gap of \$730.}
	\label{fig:economic-near-miss}
\end{figure}

\vspace{0.3cm}

\begin{econbox}[The regime signal moved in the right direction, with the inherent one-day lag.]
	
	A treasury or trader watching
	$P(\text{stress})$ climb from under $1\,\%$ to $3\,\%$ between the
	09-15 and 09-16 thresholds would have received a first, quantified
	early-warning signal of elevated risk on the very evening of
	the FOMC day, usable to adjust the hedging posture for the
	following days --- but not helpful for having avoided the
	shock itself. Even at $3{,}0\,\%$, the model remained
	"calm" at $97\,\%$, so this is a directional signal,
	not a level signal. This is the practical meaning of
	"reactive, not predictive": the tool says the ground has
	shifted once it has shifted, not before.
	
	\textbf{The single miss falls on the day the theoretical
		framework would have flagged as structurally unpredictable ---
		consistent with it, not further proof.} Every earlier
	section of this project (weak-form market efficiency,
	Meese \& Rogoff) predicts that a model built from past prices
	cannot see a future central bank decision coming. The
	FOMC statement was released on September 16 at 2:00~PM
	(Eastern time), the day of this chain's only Close-based
	breach. This is what the theoretical framework leads one to expect, but
	it rests on a single observation and on a breach of only
	$\approx 0{,}064\,\%$, within the typical range of gaps between the
	two price sources.
\end{econbox}

\subsection{What this means in practice for a model user}

\begin{econbox}
	A user with data up to a given Friday can compute a genuine,
	model-based price range for the next trading day, at their
	chosen confidence level, reflecting the current
	regime-weighted volatility --- but gets absolutely no
	information on whether the price will be higher or lower
	than today. The range answers "how far might it move",
	never "which way". This is not a simplification that loses
	something important --- it is the faithful summary of what has been
	validated (repeatedly, independently, and now empirically
	in this very test) and what has not.
\end{econbox}

\newpage

% ═════════════════════════════════════════════════════════════════
\section{Out-of-sample robustness --- the COVID-19 shock \\(January--April 2020)}
\label{sec:covid}
% ═════════════════════════════════════════════════════════════════

\subsection{Objective}
Where the rolling test (section~\ref{sec:rolling-test}, $n=4$,
September 2026) could not distinguish "expected
behavior with frozen parameters" from "a genuine structural
responsiveness problem", this test was designed specifically to
answer that question, using a genuinely out-of-sample model
and a window wide enough (84 trading days) to compute a
genuine empirical coverage rate.

\subsection{Why a separate, re-estimated model was needed}

The project's main MS-GARCH model (\texttt{ms\_garch.R}) was
calibrated on the full 2010--2026 sample --- including the
COVID period itself. Testing this model "on" 2020 would be
close to circular: the model's likelihood was partly shaped by
the very event being tested. A genuine robustness test requires a
model that never saw the shock.

\textbf{Fix:} \texttt{ms\_garch\_pre2020.R} re-estimates a
separate MS-GARCH($K=2$) using only data up to
2019-12-31 (2,501 observations). This model is then
\textbf{frozen} and used, unchanged, throughout the
January--April 2020 test --- exactly the same frozen-parameter
design as September's rolling test.

\vspace{0.3cm}

\begin{limitbox}
	\textbf{A known limitation of the pre-2020 fit, stated
		honestly...}
	
	The pre-2020 fit's stress-regime tail parameter is
	poorly identified: $\nu_2 = 71{,}73$, std.\ err.\ $=202{,}05$
	--- the uncertainty is nearly 3 times the estimate itself
	($p=0{,}361$), much less precise than the already weak
	$\nu_2 = 22{,}0$ (std.\ err.\ $=16{,}85$) of the full-sample
	model. Such a high $\nu$ is close to Normal tails. This suggests
	the pre-2020 decade contains too few genuine extreme
	events for the model to have learned how thick the stress
	regime's tail actually should be. This limitation is
	\textbf{not corrected} --- artificially fixing it would defeat the
	very purpose of testing what a model built from
	pre-2020 information alone would actually have produced.
	
	Otherwise, the fit is structurally consistent with the
	full-sample model: the long-run volatilities per
	regime ($\sigma_1 \approx 0{,}215$, $\sigma_2 \approx 0{,}726$) are
	close to the full-sample values ($0{,}206$, $0{,}730$),
	and the stable probabilities ($68{,}7\,\%$/$31{,}3\,\%$ vs
	$71{,}7\,\%$/$28{,}3\,\%$) are similar.
\end{limitbox}

\subsection{Test design}

\textbf{Window:} 2020-01-01 to 2020-04-30 (84 trading days) ---
starts exactly at the pre-2020 cutoff (no gap, no overlap)
and deliberately begins before the shock (January was calm), so
that the test covers both a stable period and the acute
crisis.
\textbf{Single data source:} unlike the September test,
this window falls entirely within the project's original FRED
sample --- no source-switch caveat applies to the
forecasting pipeline itself. 

%\textbf{Frozen parameters, a single \texttt{set.seed(42)} before the loop} --- same design as \texttt{rolling\_test\_chain.R}.

\subsection{Empirical coverage}

\begin{center}
	\begin{tabular}{ll}
		\toprule
		Coverage & $77/84 = 91{,}7\,\%$ (nominal target: $95\,\%$) \\
		Breaches & 7 (expected under nominal $95\,\%$: $\approx 4{,}2$) \\
		\bottomrule
	\end{tabular}	
	\end{center}
	
	\vspace{0.2cm}
	
	\textbf{The correct statistical test here is one-sided}, since the
	concern specifically involves under-coverage (too many
	breaches), not a gap in either direction:
	\[
	\text{One-sided binomial test } (H_1: \text{actual rate} > 5\,\%)~:
	p = 0{,}127 \quad (\text{two-sided}~: p = 0{,}200)
	\]
	
	\textbf{Interpretation, stated precisely:} even with the more
	favorable one-sided test, $p=0{,}127$ is above conventional
	thresholds ($0{,}05$ or even $0{,}10$) --- the data do not provide
	statistically significant evidence of under-coverage.
	\textbf{This is not a validation of good calibration} --- with
	$n=84$, the test has limited power to distinguish a true breach
	rate of $5\,\%$ from a true rate of $8\,\%$; the absence of
	significant evidence of a problem is not evidence of the absence of a
	problem. What this does say, in relation to the open question of
	section~\ref{sec:rolling-regime-reactivity}: the model did
	\textbf{not} collapse or show a dramatic coverage failure
	(e.g.\ $60$--$70\,\%$) during a genuine major crisis --- a
	reassuring result, without being decisive.
	
	\begin{figure}[H]
		\centering
		\includegraphics[width=1\linewidth]{01_covid_ribbon_ohlc_84days.png}
		\caption{Confidence bands and actual OHLC, 84 days (January--April 2020).}
		\label{fig:covid-ribbon-ohlc}
	\end{figure}
	
	\vspace{0.3cm}
	
	\subsection{Close-based breaches: two distinct groups, exactly as expected}
	
	\begin{table}[H]
		\centering
		\small
		\caption{Close-based breaches}
		\begin{tabular}{lrlrr}
			\toprule
			Date & Actual (Close) & Direction & Size & $P(\text{stress})$ at forecast \\
			\midrule
			2020-02-21 & 1.0855 & High & $0{,}02\,\%$ & $0{,}3\,\%$ \\
			2020-02-27 & 1.0977 & High & $0{,}31\,\%$ & $0{,}3\,\%$ \\
			2020-03-02 & 1.1164 & High & $0{,}90\,\%$ & $2{,}5\,\%$ \\
			2020-03-06 & 1.1319 & High & $0{,}09\,\%$ & $49{,}3\,\%$ \\
			2020-03-12 & 1.1081 & Low & $0{,}63\,\%$ & $98{,}7\,\%$ \\
			2020-03-17 & 1.0971 & Low & $0{,}30\,\%$ & $99{,}2\,\%$ \\
			2020-03-26 & 1.1025 & High & $0{,}44\,\%$ & $98{,}4\,\%$ \\
			\bottomrule
		\end{tabular}
	\end{table}
	
	\vspace{0.3cm}
	
	\textbf{Group~A (regime not yet reacted, $P(\text{stress})<50\,\%$):}
	02/21, 02/27, 03/02, 03/06 (borderline, $49{,}3\,\%$) --- consistent with the
	reactive-not-predictive lag already documented
	(section~\ref{sec:rolling-test}).
	
	\vspace{0.2cm}
	
	\textbf{Group~B (regime already flagged as high stress,
		$P(\text{stress})>95\,\%$):} 03/12, 03/17, 03/26 --- the model
	correctly identified acute stress \textbf{and yet} produced
	a range too narrow for the realized move.
	
	
	
	\begin{figure}[H]
		\centering
		\includegraphics[width=1\linewidth]{02_breach_groups_A_B.png}
		\caption{The two breach groups --- Group~A (regime not yet reacted) vs Group~B (regime already flagged, range still too narrow).}
		\label{fig:breach-groups}
	\end{figure}
	
	\vspace{0.3cm}
	
	\subsection{Beyond the Close: intraday High~/~Low reveals a much wider gap}
	
	Actual OHLC data (Pound Sterling Live --- FRED has no
	intraday High/Low) were obtained to check the forecast
	bands against the full daily range, not just the Close.
	
	\begin{center}
		\begin{tabular}{lr}
			\toprule
			Close-based breaches & 7/84 ($8{,}3\,\%$) \\
			Days with High $>$ upper bound & 12/84 ($14{,}3\,\%$) \\
			Days with Low $<$ lower bound & 7/84 ($8{,}3\,\%$) \\
			UNIQUE days touching outside the band (Close OR High OR Low) & \textbf{19/84 ($22{,}6\,\%$)} \\
			\bottomrule
		\end{tabular}
	\end{center}
	
	\vspace{0.2cm}
	
	\textbf{Nearly one day in four touched outside the $95\,\%$
		band at some intraday point --- about $2{,}7$ times the
		Close-based breach rate.} This is a substantially wider gap
	than the smaller September test sample suggested, and
	reinforces the already established finding (section~\ref{sec:price-range},
	section~\ref{sec:rolling-test}): the model is built for, and should
	only be judged on, close-to-close moves; intraday
	reality is noticeably more volatile than a Close-only
	coverage check would suggest.
	
	Most of the additional breaches visible only via
	High/Low cluster tightly around known crisis dates ---
	days where the Close returned inside the band by the end of the
	session, but where the price touched outside the band during the
	day.
	
	\begin{limitbox}
		\textbf{A data-sourcing caveat, stated precisely:}
		2020-03-19 --- the most extreme day of the window (PSL
		Low $=1{,}0653$, the crisis trough) --- shows a
		Low breach but \emph{not} a Close breach, implying
		that FRED's own value for that date differed enough
		from PSL's Close ($1{,}0656$) to stay inside
		the band. This is consistent with, and possibly amplified during,
		the FRED/PSL gap already documented (section~\ref{sec:price-range}) ---
		cross-source noise may matter more on the most
		volatile days, one more reason why the Close-based
		coverage figure (single-source, FRED only) is the more
		reliable of the two views for judging the model itself.
	\end{limitbox}
	
	\vspace{0.3cm}
	
	\begin{figure}[H]
		\centering
		\includegraphics[width=1\linewidth]{03_breach_rate_by_definition.png}
		\caption{Breach rate by definition --- $8{,}3\,\%$ / $14{,}3\,\%$ / $8{,}3\,\%$ / $22{,}6\,\%$.}
		\label{fig:breach-rate-by-definition}
	\end{figure}
	
	\subsection{Economic interpretation}
	
	\begin{econbox}
		\begin{itemize}
			\item \textbf{A treasury relying on this model throughout the
				crisis would have been "surprised" close to close
				about 1 day in 12} (7/84), close to but somewhat above
			the 1-in-20 promised by the model --- a modest, not
			spectacular, gap. But if that same treasury managed
			intraday risk (e.g.\ \emph{stop-loss} orders, margin
			calls triggered by any touch, not just the day's final
			print), \textbf{it would have been surprised about 1 day in 4}
			(19/84) --- a materially different risk picture depending on
			the question actually being asked of the model.
			\item \textbf{The regime signal was informative but not
				sufficient on its own on the most extreme days.} The
			Group~B cases show that even "the model correctly
			says we are $98\,\%$ certain to be in the stress
			regime" did not guarantee the band was wide enough on
			the very worst days (mid- to late March) --- a limitation
			worth knowing before treating
			"$P(\text{stress})$ is high" as a complete risk signal
			on its own.
			\item \textbf{The most useful practical takeaway}: this
			model --- like any model built purely from
			a comparatively calm decade of history --- probably underestimates
			how extreme a genuinely unprecedented crisis can become,
			even once it correctly recognizes that a crisis is
			underway. This is a point of humility about tail risk,
			more than a flaw specific to this implementation.
		\end{itemize}
	\end{econbox}
	
	\subsection{What remains open}
	
	\begin{itemize}
		\item The $\nu_2$ hypothesis (Group~B) --- not tested.
		\item The full-scale backtest over the entire 2010--2026
		sample (not just a crisis window) --- still
		the major next step planned, now carried out
		(section~\ref{sec:backtest}).
	\end{itemize}
	
	\newpage
	
	% ═════════════════════════════════════════════════════════════════
	\section{Full expanding-window backtest (2013--2026)}
	\label{sec:backtest}
	% ═════════════════════════════════════════════════════════════════
	
	This is the project's first test with genuine statistical
	power --- 3,424 one-day-ahead forecasts over 14 years --- and the first
	result allowing a real, well-founded claim about the
	model's calibration, rather than a demonstration or a
	suggestive but underpowered check.
	
	\subsection{Design}
	
	\textbf{Recursive out-of-sample, expanding-window backtest}
	\citep{west1996}: a standard, established method, not an
	\emph{ad hoc} invention. 14 \emph{folds}, each re-estimating a fresh
	MS-GARCH($K=2$) on all data from 2010 up to a given
	year-end, then forecasting every trading day of the following
	year with these \textbf{frozen} parameters:
	
	\begin{table}[H]
		\centering
		\small
		\caption{Backtest design (excerpt)}
		\begin{tabular}{lllr}
			\toprule
			Fold & Trained up to & Test year & Training obs. \\
			\midrule
			1 & 2012-12-31 & 2013 & 751 \\
			2 & 2013-12-31 & 2014 & 1,002 \\
			\multicolumn{4}{c}{\dots} \\
			8 & 2019-12-31 & 2020 & 2,501 \\
			\multicolumn{4}{c}{\dots} \\
			14 & 2025-12-31 & 2026 & 4,000 \\
			\bottomrule
		\end{tabular}
	\end{table}
	
	The first \emph{fold} deliberately starts with only 2 years
	(2010--2012, 751 observations) rather than waiting for a larger
	initial sample --- accepted as a deliberate design choice,
	with the expectation (confirmed below) that the early
	\emph{folds} would be estimated less precisely than later ones.
	
	\textbf{Computational scope decision:} only the $95\,\%$
	interval is computed per day (not the full multi-level
	$20/50/80/95\,\%$ set used for the COVID test), to keep
	runtime manageable across $\approx 3,500$ forecast steps --- a
	simpler visualization could be built from this
	single-level output if needed later.
	
	\textbf{Robustness engineering:} each of the $\approx 14$
	heavy re-estimations, as well as every \texttt{State()}/\texttt{Risk()}
	call inside the daily loop, is wrapped in
	\texttt{tryCatch} --- a single problematic year or day
	cannot silently corrupt or interrupt the run. Results
	are saved incrementally after each \emph{fold}.
	Automated warnings trigger for two red flags already
	identified earlier in the project: persistence within
	$0{,}001$ of the IGARCH boundary
	(section~\ref{sec:garch-stability}) and a tail parameter ($\nu$)
	whose standard error exceeds its own point estimate
	(section~\ref{sec:covid}).
	
	\subsection{Main result: this project's first genuinely powerful validation}
	
	\begin{center}
		\begin{tabular}{ll}
			\toprule
			Total forecasts & 3,424 \\
			Coverage & $3,270/3,424 = \mathbf{95{,}50\,\%}$ (nominal: $95\,\%$) \\
			One-sided binomial test ($H_1$: actual rate $>5\,\%$) & $\mathbf{p=0{,}919}$ \\
			\bottomrule
		\end{tabular}
	\end{center}
	
	\vspace{0.3cm}
	
	\textbf{This is a real, statistically significant result ---
		not a demonstration.} With $n=3,424$, the test has genuine
	power to detect a significantly elevated breach rate,
	and finds none: empirical coverage ($95{,}50\,\%$) sits
	almost exactly on the nominal target, and the one-sided test gives
	absolutely no evidence of under-coverage ($p=0{,}919$, far from
	any conventional threshold). This is the first point in the entire
	project where "the model appears well calibrated" can be
	stated with genuine statistical grounding, rather than as a
	demonstration (section~\ref{sec:rolling-test}, $n=4$) or a
	suggestive but underpowered check (section~\ref{sec:covid}, $n=84$, $p=0{,}127$).
	
	\vspace{0.5cm}
	
	\begin{figure}[H]
		\centering
		\includegraphics[width=1\linewidth]{02_binomial_test_full.png}
		\caption{Binomial test over the 3,424 forecasts --- the entire project's single most important result.}
		\label{fig:binomial-test-full}
	\end{figure}
	
	\newpage
	
	\subsubsection{Independent reproducibility check}
	
	Fold~8 (trained up to 2019-12-31) uses a training cutoff
	identical to the separate model built in
	section~\ref{sec:covid}. The two fits match
	\textbf{exactly}: $\alpha+\beta = 0{,}9978/0{,}9881$,
	$\nu = 7{,}10/71{,}73$, in both. This is not expected by
	construction (they were run in different scripts, at
	different times) --- it confirms that the pipeline is genuinely
	reproducible.
	
	\subsubsection{Reconciliation with the different result of section~\ref{sec:covid}}
	
	That section reported $91{,}7\,\%$ coverage (7/84 breaches)
	for January--April 2020 specifically. This backtest's
	2020 fold covers \textbf{the full calendar year} (250 days) and finds
	$95{,}2\,\%$ coverage (12 breaches). This is not
	contradictory: the implied breach rate for the remaining
	May--December 2020 period is only about $3{,}0\,\%$
	(5 breaches/166 days) --- 2020's calmer second half
	dilutes the higher concentration of breaches from the
	acute January--April shock into a full-year average close
	to nominal. Both results are correct descriptions of different
	windows.
	
	\subsection{Coverage by year}
	
	\begin{table}[H]
		\centering
		\small
		\caption{Coverage by year}
		\begin{tabular}{lrl}
			\toprule
			Year & Coverage & Breaches/$n$ \\
			\midrule
			2013 & 95.2\,\% & 12/251 \\
			2014 & 98.4\,\% & 4/250 \\
			\textbf{2015} & \textbf{91.2\,\%} & \textbf{22/251} \\
			2016 & 97.6\,\% & 6/251 \\
			2017 & 94.8\,\% & 13/249 \\
			2018 & 95.6\,\% & 11/249 \\
			2019 & 95.6\,\% & 11/249 \\
			2020 & 95.2\,\% & 12/250 \\
			2021 & 95.2\,\% & 12/249 \\
			2022 & 95.2\,\% & 12/250 \\
			2023 & 96.8\,\% & 8/249 \\
			2024 & 96.8\,\% & 8/251 \\
			2025 & 94.4\,\% & 14/250 \\
			2026 (partial, through Sept.) & 94.9\,\% & 9/175 \\
			\bottomrule
		\end{tabular}
	\end{table}
	
	\begin{figure}[H]
		\centering
		\includegraphics[width=1\linewidth]{01_coverage_by_year_events.png}
		\caption{Coverage by year, with named market events.}
		\label{fig:coverage-by-year}
	\end{figure}
	
	No year shows a dramatic coverage collapse (e.g.
	$70$--$80\,\%$), even well-known turbulent years (COVID 2020,
	the 2022 monetary tightening) --- consistent with the strong
	overall result. \textbf{2015 stands out as the weakest year.}
	
	\subsubsection{2015's underperformance is directly linked to a convergence warning --- not a coincidence to ignore}
	
	Fold~3 (trained up to 2014-12-31, testing 2015) is the
	only fold, alongside itself, to trigger the IGARCH
	boundary warning: $\alpha+\beta = 0{,}9999/0{,}9998$, essentially at
	the stationarity limit (section~\ref{sec:garch-stability} had
	documented this same fragility mode earlier in the project). \textbf{This
		is a meaningful, corroborating link, not merely a
		diagnostic curiosity}: the only fold with a numerical
	warning flag during estimation is also the only year with a
	materially degraded actual coverage ($91{,}2\,\%$ against a
	range of $94$--$98\,\%$ everywhere else). This reinforces the
	finding that this script's built-in automated convergence
	warnings detect something practically relevant, not a mere
	cosmetic estimation detail.
	
	\newpage
	
	\subsection{A recurring structural finding: $\nu_2$ is persistently hard to identify}
	
	The "stress-regime tail parameter poorly identified"
	warning ($SE(\nu_2) > \nu_2$, i.e.\ a $t$-type ratio,
	$\nu_2/SE(\nu_2)$, below 1) triggered in
	\textbf{10 of the 14 folds}. Folds~1, 3, 6, and 14 did
	not trigger it.
	
	\begin{corrbox}[Why precisely these four folds? Investigated, not assumed]
		Before writing anything about a pattern, three candidate
		explanatory variables were checked against the ratio
		$\nu_2/SE(\nu_2)$ across the 14 folds: training sample
		size, Regime~2's own GARCH persistence
		($\alpha+\beta_2$), and the point estimate of $\nu_2$
		itself.
		
		\textbf{A first pass showed a strong correlation with
			Regime~2's persistence ($r=-0{,}95$)} --- but this was
		checked further rather than reported as is, and turned out to be
		\textbf{entirely driven by fold~6's status as an extreme
			outlier} ($\alpha+\beta_2 = 0{,}436$, far outside the
		$0{,}93$--$0{,}999$ range of every other fold).
		Removing fold~6 alone collapses the correlation to $r=+0{,}22$
		(weak, and the sign flips) --- confirming that this apparent
		pattern was a single-point artifact, not a real relationship. Sample
		size showed essentially no correlation either
		($r=-0{,}20$).
		
		\textbf{The more precise --- and more sobering ---
			reframing:} since $\nu_2/SE(\nu_2)$ is mathematically a $t$
		statistic for testing $\nu_2=0$, computing approximate
		$p$-values for the 14 folds shows that \textbf{only 2 of 14 (folds~3
			and~6) reach conventional statistical significance
			($p<0{,}05$)}. The two folds that "passed" the
		script's warning threshold besides these (folds~1 and~14)
		are, in fact, also not statistically well identified
		($p \approx 0{,}16$ and $p \approx 0{,}19$) --- they simply
		crossed the script's own threshold (ratio $>1$, equivalent to a
		weak bar of $p \approx 0{,}32$), which is itself not a
		meaningful threshold for genuine statistical confidence.
		
		\textbf{Honest conclusion:} $\nu_2$ is essentially
		unidentifiable across almost all 14 folds, not
		simply "hard to identify in most folds while being
		correct in a few". No single explanatory variable
		among those checked accounts for which folds sit marginally
		above or below the (arbitrary) warning threshold --- this
		reads better as noise around a weak boundary, not a
		structural distinction between "good" and "bad"
		folds. The practical finding from the previous
		section~\ref{sec:backtest} nonetheless stands:
		the model's \emph{average} one-step coverage is excellent despite
		this --- but confidence in exactly how it prices the
		stress regime's most extreme outcomes should remain low
		throughout the sample, not just in a minority of
		folds.
	\end{corrbox}
	
	\begin{figure}[H]
		\centering
		\includegraphics[width=1\linewidth]{04_nu2_precision_by_fold.png}
		\caption{$\nu_2$ precision by fold --- essentially unidentifiable across the sample.}
		\label{fig:nu2-precision}
	\end{figure}
	
	\vspace{0.3cm}
	
	\subsection{The fold~6 anomaly (2018)}
	
	Fold~6's Regime~2 (trained on 2010--2017) showed striking
	behavior, different from every other fold:
	$\alpha+\beta_2 = 0{,}436$ (vs $0{,}93$--$0{,}999$ elsewhere) --- very
	low internal GARCH persistence --- while $\nu_2=99{,}95$
	was, unusually, \textbf{well} identified ($SE=10{,}08$, smaller
	than the estimate, the opposite of the previous section's
	recurring pattern). This was resolved once the transition matrix
	was logged
	
	
	\begin{center}
		\small
		\begin{tabular}{|p{5cm}|p{5cm}|}
			\hline
			$P(1\to2)$ & $0{,}0020$ (rare: an entry into regime~2 expected every $\approx 500$ trading days, roughly every 2 years) \\
			\hline
			$P(\text{stay in }2)$ & $0{,}9960$ (very sticky: expected duration once entered $\approx 250$ days, about 1 year) \\
			\hline
		\end{tabular}
	\end{center}
	
	\vspace{0.4cm}
	
	With only 249 test days in 2018 and an entry probability of
	$0{,}002$/day, the \emph{expected} number of entries into
	regime~2 during that specific year is only \textbf{0.5} ---
	meaning it was entirely plausible, by chance alone, that
	no full entry would occur in 2018. This matches exactly what was
	observed (regime~2 probability never exceeding $25{,}6\,\%$,
	never fully switching over).
	
	\vspace{0.4cm}
	
	\begin{figure}[H]
		\centering
		\includegraphics[width=1\linewidth]{05_fold6_regime2_outlier.png}
		\caption{Fold~6's Regime~2 --- a resolved structural edge case, not a model flaw.}
		\label{fig:fold6-outlier}
	\end{figure}
	
	\vspace{0.3cm}
	
	\textbf{What this reveals about fold~6's Regime~2, now with
		a coherent story}: combined with its low internal GARCH
	persistence (shocks fade almost immediately) and its high,
	well-identified $\nu$ (near-Normal), this fold's Regime~2 does
	not resemble the "acute stress" state found in most
	other folds --- it resembles a \textbf{rare,
		structurally distinct, but calm and stable} state once entered, not a
	panic regime. This is a legitimate, different
	characterization that the 2010--2017 training window happened to
	support, not a model flaw. Nor did this translate into
	any visible degradation of the 2018 forecasts ($95{,}6\,\%$
	coverage, $1{,}75\,\%$ average band width, both entirely
	typical) --- consistent with a regime that, that year, simply never
	got triggered rather than one that malfunctioned.
	
	\newpage
	
	\begin{limitbox}
		\textbf{Process note, stated plainly:} this section was
		first drafted with placeholder language
		("pending a re-run") before the corrected script had
		actually been run --- a sequencing error flagged directly during
		this project's own review process. The lesson
		learned: documentation edits that depend on
		results not yet generated should wait for those results, not
		anticipate them, even when the fix itself is already coded.
	\end{limitbox}
	
	\subsection{Economic interpretation --- for a reader checking these results against actual market history}
	
	\begin{econbox} The headline figure, translated
		A treasury, trader, or risk manager running
		this model over the last 14 years would have seen it keep
		almost exactly its stated promise: about 1 surprise per 20 trading
		days, matching the announced $95\,\%$ confidence
		level. This is a materially stronger claim than anything
		established earlier in the project --- no longer "the
		model looks reasonable" but "the model's stated
		uncertainty has proven empirically accurate across
		3,424 independent daily tests".
	\end{econbox}
	
	\vspace{0.3cm}
	
	\begin{econbox}
		The weak year (2015) coincides with real, named market stress.
		2015 was not a random weak point. It contains the
		\textbf{Swiss National Bank's sudden removal of the
			EUR/CHF floor} (January 15, 2015) --- a shock that,
		though centered on the franc, triggered widespread FX
		volatility spillovers in the euro area --- alongside the
		\textbf{launch of the ECB's full quantitative easing}
		(announced January 2015) and the acute phase of the
		\textbf{Greek sovereign debt crisis},
		culminating in capital controls and the July 2015
		referendum. Three distinct, significant euro-area shocks
		concentrated in a single year constitute a plausible, concrete
		explanation for why this particular fold's underlying
		volatility dynamics were harder to capture
		(the IGARCH boundary warning from the previous section) and
		why realized coverage dropped to $91{,}2\,\%$ --- not a
		mysterious statistical fluke, but a genuinely unusual year by
		the standards of the surrounding decade.
	\end{econbox}
	
	\newpage
	
	\begin{econbox}Other named events the model absorbed without breaking
		Several other major, well-documented,
		market-moving events fall in years showing
		\textbf{no degraded coverage} --- worth naming explicitly since it
		reinforces the main finding:
		\begin{itemize}
			\item \textbf{2013} (the "\emph{taper tantrum}"): the
			Fed's May 2013 signal that QE asset purchases would
			slow triggered a sharp, global repricing of risk ---
			2013 coverage ($95{,}2\,\%$) remained close to
			nominal.
			\item \textbf{2016} (the Brexit referendum, June 23, and the
			U.S. presidential election, November): two major, sudden
			\emph{risk-off}/\emph{risk-on} episodes --- 2016 shows the
			\emph{best} coverage of any year in the whole sample
			($97{,}6\,\%$).
			\item \textbf{2018} (the Italian political crisis, May 2018;
			rising U.S.--China trade tensions): coverage
			($95{,}6\,\%$) unaffected.
			\item \textbf{2020} (COVID-19): already examined in detail
			in section~\ref{sec:covid} --- the full-year coverage
			($95{,}2\,\%$) looks unremarkable precisely because the
			February--April acute shock is diluted by a calmer
			second half.
			\item \textbf{2022} (Russia's invasion of Ukraine,
			February 2022, and the Fed's most aggressive
			tightening cycle in decades, culminating with EUR/USD
			briefly reaching parity in September): coverage
			($95{,}2\,\%$) still landed almost exactly on
			target.
		\end{itemize}
	\end{econbox}
	
	\vspace{0.3cm}
	
	\begin{limitbox}
		\textbf{A recency caveat, stated plainly:} for 2023
		and beyond, this \\document does not name specific market
		events with the same confidence as the years above ---
		these years fall within or near the period where claims
		should be independently verified rather than
		asserted from memory. The \emph{statistical} results for
		these years ($2023$: $96{,}8\,\%$, $2024$: $96{,}8\,\%$, $2025$: $94{,}4\,\%$,
		partial 2026: $94{,}9\,\%$) are read directly from this
		project's own backtest output and are not in question; only the
		specific factual narrative attached to them is treated
		here with more caution.
	\end{limitbox}
	
	\newpage
	
	\begin{econbox}[Different readers, different practical takeaways]
		\begin{itemize}
			\item \textbf{A corporate treasury hedging a recurring
				EUR/USD exposure} can read this backtest as evidence that
			sizing hedges on this model's $95\,\%$ band, updated
			daily, would have left them "surprised" on
			the wrong side about as often as advertised --- about
			once a year, roughly $12$--$14$ times, not dramatically more.
			\item \textbf{A risk manager setting VaR-type limits}
			gets a stronger result: the one-sided binomial test
			means there is no statistical basis, across
			14 years including two genuine \\global-scale crises, for
			claiming that this model's stated $95\,\%$ threshold
			systematically underestimates risk --- a claim many
			simpler volatility models cannot support once
			actually tested this way.
			\item \textbf{A trader seeking a directional edge} gets
			nothing new here either, in either direction --- this
			backtest, like every other step of this project, says nothing about
			direction (sections~\ref{sec:baseline},~\ref{sec:controls}). A
			well-calibrated range is not, and never claimed to
			be, a trading signal.
			\item \textbf{A researcher or reviewer} should weigh the
			strong headline result against the previous section's honest
			caveat: the mechanism generating the extreme-tail
			forecasts ($\nu_2$) remains persistently hard to
			pin down across the sample, even though the
			\emph{average} coverage comes out correct. Good
			calibration on average is not the same claim as
			"every model component is precisely estimated" ---
			both are true here, and this document deliberately avoids
			letting the strong headline figure mask the
			weakest component.
		\end{itemize}
	\end{econbox}
	
	\newpage
	
	% ═════════════════════════════════════════════════════════════════
	\section{General conclusion}
	\label{sec:conclusion}
	% ═════════════════════════════════════════════════════════════════
	
	\subsection{What this project establishes}
	
	The two objectives stated in the introduction (\S\ref{sec:intro}) are
	delivered.
	
	\begin{econbox}[Objective (a) --- quantitatively measuring a market fear~/~volatility shift.] 
		
		A GARCH(1,1) with Student's~$t$ innovations
		($\hat\nu \approx 7{,}04$) is specified and validated on
		2010--2026 EUR/USD log returns (section~\ref{sec:garch}): clean
		diagnostics, fat tails correctly absorbed, decisive model
		comparison. A temporal stability test over five sub-periods
		(section~\ref{sec:garch-stability}) shows point estimates
		that vary, but which a formal likelihood-ratio test
		does not establish as statistically unstable ($p=0{,}277$) --- a
		nuance kept intact rather than resolved in a convenient
		direction.
		
		From this foundation, two distinct and complementary results are
		established regarding return predictability:
		\begin{itemize}
			\item the \textbf{direction} of the one-day return is not
			predictable from the volatility shock, with or without
			control variables (sections~\ref{sec:baseline} and~\ref{sec:controls}) --- a
			robust null result, consistent with weak-form market
			efficiency and the established FX literature
			\citep{meeserogoff1983,fama1984};
			\item the \textbf{magnitude} of the one-day return \emph{is}
			predictable from the level of conditional volatility
			($R^2=0{,}090$, section~\ref{sec:magnitude}) --- an internal
			consistency result of the GARCH, not a new discovery, but
			directly exploitable.
		\end{itemize}
	\end{econbox}
	
	\vspace{0.3cm}
	
	\begin{econbox}[Objective (b) --- an operational use] 
		
		Two regime detection layers, honestly labeled "low~/~normal~/~high
		volatility"
		(\S\ref{sec:regime-hmm}) --- a three-state Gaussian HMM
		(section~\ref{sec:regime-hmm}) and a two-regime Markov-Switching
		GARCH following \citet{haasmittnikpaolella2004}, implemented via
		\citet{ardia2019} (section~\ref{sec:msgarch}) --- feed a
		genuine, computed operational price range, not an
		illustrative one (section~\ref{sec:price-range}): $[1{,}1533~;~1{,}1675]$
		as of 2026-09-11, whose width "breathes" with the regime
		(stress~/~calm width ratio of $1{,}88\times$ to $2{,}09\times$
		depending on the confidence level).
	\end{econbox}	
	
	\newpage
	
	This range is validated at three nested scales, from
	demonstration to proof:
	\begin{enumerate}
		\item a rolling chain test with 4 forecasts
		(section~\ref{sec:rolling-test}) --- an end-to-end
		demonstration, not proof of calibration, with a marginal
		miss dissected into two distinct channels rather than left as a
		single verdict;
		\item an out-of-sample robustness test on the COVID-19 crisis
		(section~\ref{sec:covid}, $n=84$): coverage of $91{,}7\,\%$,
		$p=0{,}127$ --- suggestive, not decisive, with
		acknowledged limited statistical power;
		\item a full expanding-window backtest over 14 years
		(section~\ref{sec:backtest}, $n=3,424$, the method of
		\citet{west1996}): coverage of $95{,}50\,\%$ against a
		nominal target of $95\,\%$, $p=0{,}919$ --- the project's first
		validation with genuine statistical power, \\absorbing without
		breaking the 2013 \emph{taper tantrum}, the 2016 Brexit and
		U.S. election, the 2018 Italian political crisis, the 2020
		COVID shock, and the 2022 Russia-Ukraine shock, while
		honestly flagging its one weak year (2015, $91{,}2\,\%$) and
		linking it both to an internal numerical
		warning (the IGARCH boundary) and to three real, named
		macroeconomic shocks (the Swiss National Bank, the ECB's
		QE launch, the Greek debt crisis).
	\end{enumerate}
	
	\newpage
	
	\subsection{A remark on method, not just on results}
	
	A common thread runs through this project's eleven steps,
	more visible reading it start to finish than within
	any single section: every significant result was subjected to at
	least one serious attempt to challenge it before being
	accepted --- a test configuration error corrected
	(section~\ref{sec:garch}), a poorly chosen diagnostic
	threshold corrected (section~\ref{sec:garch}), a
	plausible economic narrative ultimately disproven by
	independent evidence (section~\ref{sec:garch-stability}), an
	entire conceptual framing revised before any code was written
	(section~\ref{sec:regime-hmm}), a distortion channel quantified rather
	than assumed negligible (section~\ref{sec:rolling-test}), a
	hypothesis explicitly left unverified rather than presented as settled
	(section~\ref{sec:covid}). The fact that the final backtest's
	headline result ($p=0{,}919$) comes after this accumulation of
	checks --- not instead of them --- is what gives it its weight.
	
	\subsection{What remains open}
	
	Consolidated from earlier sections, with no attempt to
	present it as thinner than it is:
	
	\begin{itemize}
		\item The mechanism behind $\nu_2$'s persistent
		imprecision (the stress-regime tail parameter) is not resolved ---
		flagged consistently in sections~\ref{sec:msgarch},
		\ref{sec:covid}, and~\ref{sec:backtest}, never masked by the
		otherwise excellent average coverage result.
		.
		\item The Viterbi path used for the annual regime
		counts (section~\ref{sec:msgarch}) is not produced by any
		script currently under version control --- a reproducibility
		gap flagged explicitly rather than hidden, on the same footing
		as the missing saved standard errors for the stability test
		(section~\ref{sec:garch-stability}).
		\item No independent (non-FRED) cross-check has been
		conducted across the full 14 years of the complete backtest --- a
		deliberate scope choice, not an oversight, but one that leaves a
		spot-check against a handful of well-known extreme dates as a
		cheaper alternative, not yet carried out
		(section~\ref{sec:backtest}).
		\item The hypothesis that an insufficiently thick pre-2020
		$\nu_2$ explains the COVID Group~B breaches
		(section~\ref{sec:covid}) remains untested.
	\end{itemize}
	
	\vspace{0.3cm}
	
	\textbf{What this project does not claim to be:} a trading signal.
	No step, from start to finish, produced or sought to produce an
	exploitable directional edge --- the robust null result on
	direction (sections~\ref{sec:baseline} and~\ref{sec:controls}) and the
	explicit reminder at every later step that the price range carries
	no directional information are the constant confirmation of this. What
	this project delivers is narrower and, on the basis of the
	$n=3,424$ backtest, better established: \textbf{\textit{a quantitative,
			testable measure of market volatility and regime, and a price
			range whose calibration has been empirically verified,
			not merely assumed.}}
	
	% ═════════════════════════════════════════════════════════════════
	\newpage
	\section*{Appendix --- supplementary figures}
	\addcontentsline{toc}{section}{Appendix --- supplementary figures}
	\renewcommand{\thefigure}{A.\arabic{figure}}
	\setcounter{figure}{0}
	% ═════════════════════════════════════════════════════════════════
	
	This appendix collects figures produced over the course of the
	project but not essential to reading the body text, following
	the \emph{Graphs\_Ways.md} arbitration: a figure appears here only if it
	supports a specific claim cited in the body --- a
	skeptical reader must be able to find it. Purely
	diagnostic figures, with no direct link to a claim in the
	text (useful for reproducibility, not for reading),
	stay outside this document.
	
	
	
	\subsection*{GARCH(1,1) (\S\ref{sec:garch})}
	
	\begin{figure}[H]
		\centering
		\includegraphics[width=1\linewidth]{04_standardized_residuals_studentt.png}
		\caption{Standardized residuals of the Student's~$t$ model.}
	\end{figure}
	
	\vspace{0.5cm}	
	
	\begin{figure}[H]
		\centering
		\includegraphics[width=1\linewidth]{07_model_comparison_aic_bic.png}
		\caption{AIC/BIC comparison, Normal vs Student's~$t$.}
	\end{figure}
	
	\vspace{0.5cm}	
	
	\begin{figure}[H]
		\centering
		\includegraphics[width=1\linewidth]{08_kurtosis_comparison.png}
		\caption{Comparison of empirical and theoretical kurtosis.}
	\end{figure}
	
	\vspace{0.5cm}	
	
	\subsection*{Baseline regression (\S\ref{sec:baseline})}
	
	\begin{figure}[H]
		\centering
		\includegraphics[width=1\linewidth]{05_distribution_deltay_percentiles.png}
		\caption{Distribution of $\Delta Y(t)$ and percentiles used for the exercise in \S\ref{sec:baseline} (economic effect at $\pm0{,}05$/$\pm0{,}20$).}
	\end{figure}
	
	\vspace{0.5cm}	
	
	\subsection*{Regression with controls (\S\ref{sec:controls})}
	
	\begin{figure}[H]
		\centering
		\includegraphics[width=1\linewidth]{02_aic_bic_comparison.png}
		\caption{AIC/BIC comparison, re-estimated baseline against the extended specification with controls.}
	\end{figure}
	
	\vspace{0.5cm}		
	
	\begin{figure}[H]
		\centering
		\includegraphics[width=1\linewidth]{03_forest_plot_compact.png}
		\caption{Compact \emph{forest plot} of the extended specification's coefficients.}
	\end{figure}
	
	\vspace{0.5cm}	
	
	\subsection*{Magnitude regression (\S\ref{sec:magnitude})}
	
	\begin{figure}[H]
		\centering
		\includegraphics[width=1\linewidth]{01_specA_scatter_deltay.png}
		\caption{Spec~A --- scatter plot of $|r(t+1)|$ against $\Delta Y(t)$ (diluted relationship, \S\ref{sec:magnitude}).}
	\end{figure}
	
	\vspace{0.5cm}	
	
	\begin{figure}[H]
		\centering
		\includegraphics[width=1\linewidth]{05_predicted_magnitude_percentiles.png}
		\caption{Predicted magnitude at the $\sigma(t)$ percentiles used in the worked example of \S\ref{sec:magnitude}.}
	\end{figure}
	
	\vspace{0.5cm}		
	
	\subsection*{HMM regime detection (\S\ref{sec:regime-hmm})}
	
	\begin{figure}[H]
		\centering
		\includegraphics[width=1\linewidth]{03_density_logsigma_by_regime.png}
		\caption{Density of $\log\sigma(t)$ by regime.}
	\end{figure}
	
	\vspace{0.5cm}		
	
	\begin{figure}[H]
		\centering
		\includegraphics[width=1\linewidth]{05_transition_matrix_heatmap.png}
		\caption{HMM transition matrix, as a heatmap.}
	\end{figure}
	
	\vspace{0.5cm}		
	
	\begin{figure}[H]
		\centering
		\includegraphics[width=1\linewidth]{07_recent_period_zoom.png}
		\caption{Zoom on the regime classification of the most recent observations.}
	\end{figure}
	
	\vspace{0.5cm}		
	
	\subsection*{Markov-Switching GARCH (\S\ref{sec:msgarch})}
	
	\begin{figure}[H]
		\centering
		\includegraphics[width=1\linewidth]{01_smoothed_prob_regime2_fullsample.png}
		\caption{Smoothed probability of regime~2, full sample (unzoomed view, complement to Figure~\ref{fig:covid-detection-zoom}).}
	\end{figure}
	
	\vspace{0.5cm}		
	
	\begin{figure}[H]
		\centering
		\includegraphics[width=1\linewidth]{07_information_criteria_delta.png}
		\caption{Direct AIC/BIC gap between MS-GARCH and simple GARCH ---
			$\Delta$AIC~$=-13{,}51$, $\Delta$BIC~$=+18{,}18$ --- referenced from
			\S\ref{sec:msgarch}, replacing a discarded raw bar
			chart (the gap was visually imperceptible there at the scale of
			the AIC/BIC values themselves).}
		\label{fig:annex-info-criteria-delta}
	\end{figure}
	
	\vspace{0.5cm}		
	
	\subsection*{Operational price range (\S\ref{sec:price-range})}
	
	\begin{figure}[H]
		\centering
		\includegraphics[width=1\linewidth]{02_expost_ohlc_vs_range.png}
		\caption{\emph{Ex-post} OHLC vs the forecast range, 2026-09-14.}
	\end{figure}
	
	\vspace{0.5cm}		
	
	\begin{figure}[H]
		\centering
		\includegraphics[width=1\linewidth]{03_range95_by_scenario.png}
		\caption{$95\,\%$ range by regime scenario (mixture, certain calm, certain stress).}
	\end{figure}
	
	\vspace{0.5cm}		
	
	\subsection*{Rolling chain test (\S\ref{sec:rolling-test})}
	
	\begin{figure}[H]
		\centering
		\includegraphics[width=1\linewidth]{01_rolling_bands_and_candles.png}
		\caption{Forecast bands and OHLC candles, preliminary view (variant preceding Figure~\ref{fig:range-vs-actual-chain}).}
	\end{figure}
	
	\vspace{0.5cm}		
	
	\begin{figure}[H]
		\centering
		\includegraphics[width=1\linewidth]{03_ohlc_vs_range_chain.png}
		\caption{Full OHLC (Open/High/Low/Close) vs the forecast range, chained across the 4 forecasts.}
	\end{figure}
	
	\vspace{0.5cm}		
	
	\begin{figure}[H]
		\centering
		\includegraphics[width=1\linewidth]{04_prob_stress_vs_range_width.png}
		\caption{$P(\text{stress})$ and range width, on actual scales --- a precise complement to Figure~\ref{fig:stress-prob-vs-width} (which indexes both series to base 100 for visual readability).}
	\end{figure}
	
	\vspace{0.5cm}		
	
	\subsection*{Out-of-sample robustness --- COVID-19 (\S\ref{sec:covid})}
	
	\begin{figure}[H]
		\centering
		\includegraphics[width=1\linewidth]{05_user_surprise_rate.png}
		\caption{"Surprise" rate for a model user, close vs intraday --- 1 day in 12 vs 1 day in 4 (\S\ref{sec:covid}).}
	\end{figure}
	
	\vspace{0.5cm}		
	
	\subsection*{Temporal stability of the GARCH (\S\ref{sec:garch-stability})}
	
	\begin{figure}[H]
		\centering
		\includegraphics[width=1\linewidth]{02_parameters_by_period.png}
		\caption{Parameters $\omega,\alpha,\beta,\nu$ by period.}
	\end{figure}
	
	\vspace{0.5cm}		
	
	\begin{figure}[H]
		\centering
		\includegraphics[width=1\linewidth]{03_halflife_vs_persistence.png}
		\caption{Half-life vs persistence ($\alpha+\beta$), by period.}
	\end{figure}
	
	\vspace{0.5cm}		
	
	\begin{figure}[H]
		\centering
		\includegraphics[width=1\linewidth]{04_omega_by_period.png}
		\caption{$\omega$ by period --- illustrates Period~2's imprecision.}
	\end{figure}
	
	\vspace{0.5cm}		
	
	\subsection*{Full backtest (\S\ref{sec:backtest})}
	
	\begin{figure}[H]
		\centering
		\includegraphics[width=1\linewidth]{03_alphabeta_by_fold.png}
		\caption{$\alpha+\beta$ by expanding-window backtest \emph{fold}.}
	\end{figure}
	
	% ═════════════════════════════════════════════════════════════════
	\newpage
	\section*{References}
	\addcontentsline{toc}{section}{References}
	% ═════════════════════════════════════════════════════════════════
	
	\begin{thebibliography}{99}
		
		\bibitem[Ardia et al.(2019)]{ardia2019}
		Ardia, D., Bluteau, K., Boudt, K., Catania, L., \& Trottier, D.-A.
		(2019). Markov-Switching GARCH Models in R: The MSGARCH Package.
		\textit{Journal of Statistical Software}, 91(4), 1--38.
		
		\bibitem{fama1984}
		Fama, E.~F. (1984). Forward and Spot Exchange Rates.
		\textit{Journal of Monetary Economics}, 14(3), 319--338.
		
		\bibitem{grangernewbold1974}
		Granger, C.~W.~J., \& Newbold, P. (1974). Spurious Regressions in
		Econometrics. \textit{Journal of Econometrics}, 2(2), 111--120.
		
		\bibitem[Gray(1996)]{gray1996}
		Gray, S.~F. (1996). Modeling the Conditional Distribution of Interest
		Rates as a Regime-Switching Process. \textit{Journal of Financial
			Economics}, 42(1), 27--62.
		
		\bibitem[Haas et al.(2004)]{haasmittnikpaolella2004}
		Haas, M., Mittnik, S., \& Paolella, M.~S. (2004a). A New Approach to
		Markov-Switching GARCH Models. \textit{Journal of Financial
			Econometrics}, 2(4), 493--530.
		
		\bibitem[Hamilton and Susmel(1994)]{hamiltonsusmel1994}
		Hamilton, J.~D., \& Susmel, R. (1994). Autoregressive Conditional
		Heteroskedasticity and Changes in Regime. \textit{Journal of
			Econometrics}, 64(1--2), 307--333.
		
		\bibitem[Klaassen(2002)]{klaassen2002}
		Klaassen, F. (2002). Improving GARCH Volatility Forecasts with
		Regime-Switching GARCH. \textit{Empirical Economics}, 27(2), 363--394.
		
		\bibitem[Marcucci(2005)]{marcucci2005}
		Marcucci, J. (2005). Forecasting Stock Market Volatility with
		Regime-Switching GARCH Models. \textit{Studies in Nonlinear Dynamics \&
			Econometrics}, 9(4), Article~6.
		
		\bibitem{meeserogoff1983}
		Meese, R.~A., \& Rogoff, K. (1983). Empirical Exchange Rate Models of
		the Seventies: Do They Fit Out of Sample? \textit{Journal of
			International Economics}, 14(1--2), 3--24.
		
		\bibitem[West(1996)]{west1996}
		West, K.~D. (1996). Asymptotic Inference about Predictive Ability.
		\textit{Econometrica}, 64(5), 1067--1084.
		
	\end{thebibliography}
	
	
\end{document}